\documentclass[11pt,a4paper]{article}
\usepackage[margin=27mm,headheight=15pt]{geometry}
\usepackage{fontspec}
\setmainfont{TeX Gyre Pagella}
\setsansfont{TeX Gyre Heros}
\setmonofont[Scale=0.84]{Latin Modern Mono}
\usepackage{amsmath,amssymb,amsthm,mathtools}
\usepackage{booktabs,longtable,array,graphicx,microtype,enumitem}
\usepackage[round,authoryear]{natbib}
\usepackage{xurl}
\usepackage[colorlinks=true,allcolors=blue!45!black]{hyperref}
\usepackage{xcolor}
\usepackage{fancyhdr}
\pagestyle{fancy}
\fancyhf{}
\fancyhead[L]{\small Relative-position surfaces}
\fancyhead[R]{\small RESEARCH\_SA11}
\fancyfoot[C]{\thepage}
\setlength{\parskip}{3pt plus 1pt}
\setlength{\emergencystretch}{2em}
\setlist{topsep=3pt,itemsep=2pt,parsep=0pt}
\newenvironment{resultenv}[3]{\par\medskip\phantomsection\label{#1}\noindent\textbf{Statement #1 --- #3.}\quad\textsf{[#2]}\par\smallskip\noindent\ignorespaces}{\par\medskip}
\hypersetup{pdftitle={Relative-position surfaces of pseudoline arrangements},pdfauthor={GPT-6 Astra, OpenAI, ChatGPT Pro mode}}
\title{Relative-position surfaces of pseudoline arrangements\\[5pt]\large Exact topology, extension-space comparison, and polygon reconstruction}
\author{AI-produced by GPT-6 Astra\\OpenAI, ChatGPT Pro mode}
\date{19 November 2026}
\begin{document}
\maketitle
\begin{abstract}
We determine the surface obtained by Grothendieck's corrected disk gluing
for every finite non-pencil arrangement of distinct pseudolines in the
projective plane, allowing multiple intersections. We use the variant
$\mathcal X_2$ of cote 154, scans 5 and 7. If $c_j$ counts the ordinary
unoriented relative positions passing through exactly $j$ old vertices,
then
\[
 \mathcal X_2\cong
 \#^{\,n+1+4\sum_{k\ge1}k c_{2k+1}}\mathbb {RP}^{2}.
\]
We prove that the gluing is well defined, every closed cell embeds,
and every vertex link is a four-cycle. The extension quotient is
homotopy equivalent to $\mathbb {RP}^{2}$; its specified real blow-up
at the parallel extensions is homotopy equivalent to
$\#^{n+1}\mathbb {RP}^{2}$. The dual manuscript surface is an explicitly
described subcomplex of that blow-up, and the inclusion is a homotopy
equivalence exactly when no ordinary position passes through three old
vertices. A fully enumerated six-line arrangement has
$(s_0,s_1,s_2)=(360,720,315)$ and genus 47, refuting the proposed
universal Euler formula. We also give an exact finite image criterion
for deployed polygons, reconstruct the weighted cellular graph,
prove the disk sufficiency and cluster-reversal statements, and
separate the boundary-to-interior tethers of cote 156-4 from proper
chords. All finite counterexamples come with exact code and outputs.
\end{abstract}

\tableofcontents
\newpage
\section*{Scope, conventions, and outcome}
\addcontentsline{toc}{section}{Scope, conventions, and outcome}
We work with labelled, finite arrangements of distinct tame pseudolines
in $\mathbb {RP}^{2}$. Each pair meets at precisely one transverse
crossing. ``Nontrivial'' means that the arrangement is not a pencil,
so its oriented matroid has rank three. The word ``simple'' for the
oriented matroid excludes loops and parallel old elements; a
\emph{simple arrangement} additionally excludes triple intersections.
The latter extra condition is imposed only in the explicitly named
finite computations. Relative positions are combinatorial isotopy
classes, equivalently labelled extension data under the relative
representation comparison proved in A.3. We do not count parametrized
embedded curves or points in a metric realization space.

The construction in cote 154 has several versions. Our $\mathcal X$
always means $\mathcal X_2$: every ordinary relative position labels a
face, and exceptional sweeps across old pseudolines are included.
The capped variant $\mathcal X_0$ is different. A vertex of our primal
surface has valence four; correspondingly a face of the dual graph is
a quadrangle. The primal faces themselves have arbitrary polygonal
orders. This distinction resolves the potentially ambiguous phrase
``faces of order four'' in the supplied question.

\begin{center}\small
\begin{tabular}{p{0.16\linewidth}p{0.62\linewidth}p{0.13\linewidth}}
\toprule
Item & Outcome & Statements\\\midrule
A.1 & Surface existence and exact topology for every arrangement in scope; universal $1-n$ formula false & T1--T3, E1--E5, C.01--C.07\\
A.2 / L2 & Stable-state selector false; embedded diamond-and-band comparison proved & E6\\
A.3 / L3 & Proposed local image test false; complete image and graph reconstruction proved & P4--P6\\
A.4 & Disk sufficiency, injectivity, and cluster reversal proved, with exact degenerate qualification & P1--P3\\
L1 & Universal Euler and blow-up equivalence false; exact formula and equivalence criterion proved & E4--E5, C.04\\
L4 & Literal identification of tether and chord data false; marked two-tether correspondence proved & T4\\
\bottomrule
\end{tabular}
\end{center}

Every numbered statement has a status and a proof or an opened-source
derivation. Statement E1 uses a known theorem, while its equivariant
quotient consequence is established explicitly. The conditional
formulation E5 is applied only after T3 verifies every hypothesis.
The formulas give an effective determination: P4 enumerates the entire
finite extension set from the base chirotope, and hence supplies every
$c_j$. They do not posit that the answer depends on $n$ alone.

We distinguish the usual nonorientable genus $h=2-\chi$, the number of
projective-plane summands, from the manuscript parameter $g=1-\chi$.
Thus $h=g+1$. The finite test with $\chi=-45$ has $h=47$ and $g=46$.

The work is AI-produced by GPT-6 Astra (OpenAI, ChatGPT Pro mode).
No mathematical contribution of this run is attributed to a person,
and we claim no priority. Literature authors are cited for their
published results. The supplied AI studies and transcriptions are
witnesses to check, not authorities; manuscript references below are
cote/scan locators, not an unverified identification with handwritten
page numbers.
\section{A.1. The surface and its exact topology}
% Extension-space branch: complete proofs. AI-produced by GPT-6 Astra.
% Intended for the root paper's resultenv environment (id,status,title).
\subsection{The actual extension space, its marked points, and the cubical comparison}

Throughout this subsection, $\Sigma$ is a finite arrangement of $n$ distinct
pseudolines in $\mathbb {RP}^{2}$, any two meeting transversely at precisely one
point, and not all concurrent. Thus its oriented matroid $M$ is simple of rank
three; concurrent subfamilies are allowed. We fix orientations of the old
pseudolines only to write sign vectors. Let $\mathcal C^*(M)$ denote the set of
signed cocircuits. The extension poset $P=\mathcal E(M)$ consists of rank-three
one-element extensions by a labelled element $f$ which is not a loop. Equivalently,
these are the nonzero localization maps
\[
 \sigma:\mathcal C^*(M)\longrightarrow\{0,+,-\},\qquad
 \sigma(-Y)=-\sigma(Y).
\]
The partial order is coordinatewise, with $0<+$ and $0<-$ and with $+$ and $-$
incomparable. We write $|P|$ for the geometric realization of its order complex.
The involution $\tau(\sigma)=-\sigma$ reorients only the new element. It does not
reorient any old element and does not relabel the arrangement.

For an old element $g$, the two parallel extensions are denoted $[g^+]$ and
$[g^-]$. We call every other extension \emph{ordinary}. An ordinary extension
is represented by a pseudoline distinct from the old pseudolines. We write
$\nu(D)$ for the number of old vertices on that pseudoline, equivalently the
number of antipodal cocircuit pairs on which its localization vanishes. Set
\[
 c_j=\#\{\text{ordinary extensions modulo }\tau:\nu(D)=j\},
 \qquad F(t)=\sum_{j\geq0}c_jt^j.
\]
The old vertices form a finite set $V(\Sigma)$; $\mu(v)$ is the number of old
pseudolines through $v$, and $n_i$ is the number of vertices on the $i$th old
pseudoline. In particular $\sum_v\mu(v)=\sum_i n_i$.

\begin{resultenv}{E1}{Known}{The extension quotient}
The space $|P|$ is homotopy equivalent to $S^2$. The action $\tau$ is free and
there is a homotopy equivalence
\[
 |P|/\langle\tau\rangle\simeq\mathbb {RP}^{2}.
\]
More precisely, for every old element $g$ the map
\[
 \Pi_g:P\longrightarrow Q_g:=\{g^+,g^-\}\oplus\mathcal E(M/g)
\]
is a $\tau$-equivariant homotopy equivalence on order complexes. Here $\oplus$
means the ordinal sum: the two incomparable symbols $g^+,g^-$ lie below every
element of $\mathcal E(M/g)$. On ordinary non-$g$-parallel extensions,
$\Pi_g(\sigma)$ is restriction to cocircuits $Y$ with $Y_g=0$; on $[g^\pm]$
it is the corresponding symbol $g^\pm$.
\end{resultenv}
\begin{proof}
The sphere statement is Sturmfels and Ziegler's Theorem~1.2, p.~25, together
with their Corollary~4.5, p.~40; the displayed map and its homotopy equivalence
are their Proposition~4.2, pp.~37--38 \cite{SZ1993}. We check the
hypotheses actually needed. The contraction $M/g$ has rank two. Every
one-element extension $M\cup f$ under consideration still has rank three,
and every rank-three oriented matroid program is Euclidean, as used in their
Corollary~4.5 and the proof of Proposition~4.2. The latter proposition therefore
applies to each choice of $g$, not merely to a special realizable case.

The formula for $\Pi_g$ commutes with $\sigma\mapsto-\sigma$: this negates the
restricted localization and exchanges the two suspension symbols. The
rank-two extension poset is the face poset of a circle, with sign reversal
acting antipodally. To see the rank-two assertion directly, represent the
rank-two oriented matroid by nonzero vectors in the oriented plane, with each
parallel class represented by one direction. The new vector's direction runs
around a circle; the finitely many old signed directions are the vertices
and the intervening intervals are the edges. Parallel elements and loops of
$M/g$ do not alter this circle. Suspending gives an actual $S^2$ with the
ordinary antipodal action: in suspension coordinates $(z,t)$ it is
$(z,t)\mapsto(-z,-t)$.

The action on $|P|$ is free. Indeed, $\sigma\neq-\sigma$ for every nonzero
localization, and an invariant finite chain under an order-preserving
involution would have every member fixed. Thus it fixes no simplex setwise
and no point. Both spaces are finite CW complexes and their quotients are
covered by these free two-sheeted covers. Since $|\Pi_g|$ is a homotopy
equivalence and both total spaces are simply connected, the induced quotient
map is an isomorphism on fundamental groups (both deck groups are $C_2$) and,
by the covering-space lifting isomorphisms, on every higher homotopy group.
The CW Whitehead theorem proves that the quotient map is a homotopy
equivalence. This establishes the quotient conclusion from the actual
specified equivariant map, rather than from an unmarked sphere homotopy type.
\end{proof}

\begin{resultenv}{E2}{Proved}{An embedded regular cubical and polygonal model}
The opposite poset $P^{\mathrm{op}}$ is the face poset of a finite regular CW
complex $T(M)$ with the following cells: an ordinary extension $D$ labels a
$\nu(D)$-cube, and each parallel extension labels a polygonal two-cell.
Every closed cell embeds. The barycentric subdivision of $T(M)$ is $|P|$,
and its quotient $\overline T(M)=T(M)/\langle\tau\rangle$ has one ordinary
cube per unoriented relative position and $n$ parallel polygonal two-cells.
Consequently
\[
 \boxed{F(-1)=1-n.}
\]
\end{resultenv}
\begin{proof}
Let $\sigma$ be ordinary, with zeros at $\nu$ old vertices. Distinct such
vertices involve disjoint sets of old pseudolines: otherwise the added
pseudoline would meet one old pseudoline at two distinct points. Choose
pairwise disjoint small disks about these vertices. Resolving the added
pseudoline to either side inside any chosen disk and leaving it fixed
elsewhere preserves its single transverse intersection with every old
pseudoline. These resolutions can be made independently. At each vertex we
may leave the passage unresolved, or choose either resolution. They give
exactly $3^\nu$ extensions above $\sigma$: no sign outside its zero set may
change in a greater localization, and its signs on the zero pairs determine
the extension. The latter uniqueness is the localization determination
theorem; it is also immediate here by gluing the specified local arcs back
to the fixed complement. With order reversed, these are precisely the
nonempty faces of $[-1,1]^\nu$, where a zero coordinate means a free cube
coordinate. Thus the proper upper interval $P_{>\sigma}$ is the face poset
of the boundary of that cube, with order reversed.

For a parallel extension $[i^\pm]$, Sturmfels--Ziegler Lemma~4.1, p.~37,
identifies $P_{>[i^\pm]}$ with $\mathcal E(M/i)$, a rank-two circle. It is
therefore a polygonal boundary, and adding its cone gives a two-cell.

We make the embedding assertion explicit. In the finite simplicial complex
$|P^{\mathrm{op}}|$, for each element $x$ take the principal-ideal
subcomplex
\[
 C_x=|\Delta((P^{\mathrm{op}})_{\leq x})|.
\]
It is literally the cone, with vertex $x$, on
$|\Delta((P^{\mathrm{op}})_{<x})|$. The preceding interval calculations say
that this is the barycentrically subdivided closed cube, or the subdivided
closed polygon, appropriate to $x$. Its open cell is $C_x$ minus that
boundary. Every point of the order complex belongs to exactly one such
open cell, as seen from the largest element of its supporting chain.
Intersections of closed cells are unions of their common closed faces.
Thus these are embedded characteristic balls in a finite regular CW
structure; no assertion that an arbitrary colimit of embeddings is embedded
is being used.

For any $x$, choose a cocircuit on which its localization is nonzero. Every
extension above $x$ has that same sign. Hence $C_x$ contains no pair of
antipodal extensions and no pair of points identified by $\tau$. The quotient
map therefore remains injective on each $C_x$ and gives the asserted regular
CW structure on $\overline T(M)$. Its Euler characteristic is
$F(-1)+n$; E1 makes this equal to $\chi(\mathbb {RP}^2)=1$.
\end{proof}

\begin{resultenv}{E3}{Proved}{The derivative identity from a fixed old vertex}
For every $v\in V(\Sigma)$,
\[
 \sum_{D\ni v}(-1)^{\nu(D)-1}=\mu(v),
\]
where the sum is over ordinary unoriented relative positions through $v$.
In particular
\[
 \boxed{F'(-1)=\sum_{v\in V(\Sigma)}\mu(v)=\sum_i n_i.}
\]
\end{resultenv}
\begin{proof}
Choose an oriented cocircuit $Y$ at $v$ and an old element $g$ through $v$.
Let $P_v=\{\sigma\in P:\sigma(Y)=0\}$, with its induced order. The condition
is independent of the choice between $Y$ and $-Y$. Since $Y_g=0$, the map
$\Pi_g$ of E1 restricts to
\[
 P_v\longrightarrow\{g^+,g^-\}\oplus R_v,
 \qquad R_v=\{\rho\in\mathcal E(M/g):\rho(Y)=0\}.
\]
In rank two, $R_v$ consists of two antipodal parallel extensions and is an
antichain with two elements. Thus the target order complex is a circle.

We verify that this restriction still induces a homotopy equivalence; this
does not follow merely by restricting an arbitrary homotopy equivalence.
Sturmfels and Ziegler prove Proposition~4.2 using their Lemma~3.2, p.~29
(Babson's fiber criterion). The proof verifies that each point fiber of
$\Pi_g$ is contractible, and that
\[
 \Pi_g^{-1}(q)\cap P_{\geq p}
\]
is contractible whenever $\Pi_g(p)<q$. For a restriction to the full inverse
image of any subposet of $Q_g$, each of these fibers and slices is unchanged.
The same criterion therefore applies to the restricted map. In the present
case one can see the slices especially explicitly: the only comparisons
in the target are $g^\pm<\rho$, and the corresponding slice is the singleton
$\{\kappa_\pm(\rho)\}$ of their Lemma~4.1. We have proved $|P_v|\simeq S^1$.
The sign action is free, so $\chi(|P_v|/\langle\tau\rangle)=0$.

In the opposite induced poset, an ordinary extension through $v$ has a
$(\nu(D)-1)$-cube as closure: the zero at $v$ must remain zero and the other
zeros can be resolved independently. A parallel extension at an old line
through $v$ has a one-cell as closure, since its proper upper interval in
$P_v$ is exactly the two rank-two localizations vanishing at the cocircuit
of $v$. There are $\mu(v)$ such one-cells modulo $\tau$. The principal-ideal
construction of E2 again proves that all these cells embed. Therefore
\[
 0=\chi(|P_v|/\langle\tau\rangle)
   =\sum_{D\ni v}(-1)^{\nu(D)-1}-\mu(v).
\]
Summing over $v$ counts an ordinary extension with $\nu(D)$ zeros exactly
$\nu(D)$ times and gives the derivative identity.
\end{proof}

\begin{resultenv}{E4}{Proved}{The marked puncture and real blow-up comparison}
Choose pairwise disjoint small open two-disks about the $n$ parallel points
in $\overline T(M)$; these points lie in genuine open two-disk neighborhoods.
Let $U(M)$ be their complement. There is a homotopy equivalence of pairs
\[
 \bigl(\mathbb {RP}^2\setminus\textstyle\bigcup_{i=1}^n\mathring B_i,
       \coprod_i\partial B_i\bigr)
 \longrightarrow
 \bigl(U(M),\coprod_i\partial B_i'\bigr)
\]
whose restrictions to the boundary circles are homeomorphisms matching the
old-line labels. Replace each removed disk by a M\"obius band, attaching its
boundary by the same circle parametrization, and denote the resulting CW
complex by $B(M)$. Then
\[
 B(M)\simeq\#^{\,n+1}\mathbb {RP}^2,
 \qquad\chi(B(M))=1-n.
\]
This statement concerns a specified operation at regular parallel points;
it does not assume that the whole extension complex is a surface.
\end{resultenv}
\begin{proof}
We give the marked argument because an unmarked homotopy equivalence to
$\mathbb {RP}^2$ alone would not justify puncturing.
Work first in the oriented complex $|P|$, and fix $g$. Put
$C=\mathcal E(M/g)$, a circle with antipodal sign action, regarded as a
face poset. Its zero-dimensional cells will be denoted $\rho$ and its
one-dimensional cells $\gamma$. For a closed edge $\gamma$ define
\[
 Q_\gamma=\{g^+,g^-\}\oplus C_{\leq\gamma},\quad
 R_\gamma=|\Pi_g^{-1}(Q_\gamma)|.
\]
For a vertex $\rho$ define
\[
 Q_\rho=\{g^+,g^-,\rho\},\qquad
 A_\rho=|\Pi_g^{-1}(Q_\rho)|.
\]
Each $|Q_\gamma|$ is the suspension of an interval, hence a closed disk;
each $|Q_\rho|$ is an interval with endpoints $g^+,g^-$. By the restricted
fiber argument in E3, every $R_\gamma$ and every $A_\rho$ is contractible.
The $R_\gamma$ cover $|P|$. Adjacent sectors meet in $A_\rho$; nonadjacent
sectors meet only in the two points $[g^+],[g^-]$. These assertions hold as
literal equalities of simplicial subcomplexes, because a chain in a circle's
face poset lies over one closed edge.

We describe the local marked neighborhoods. The star of any parallel
extension $[h]$ is a cone on the rank-two circle $\mathcal E(M/h)$, by
Lemma~4.1, and $[h]$ is minimal in $P$. Thus its star is a genuine open
two-disk neighborhood in the entire order complex. For $h\neq g$, its image
under $\Pi_g$ is the rank-two parallel vertex $\rho$ corresponding to
$v=g\cap h$. In its link $\mathcal E(M/h)$, evaluation at the cocircuit of
$v$ has two zeros and separates the circle into the positive and negative
arcs. The two zeros form the link of $[h]$ in $A_\rho$; the two closed arcs
form its links in the two adjacent $R_\gamma$. Hence $[h]$ is locally an
interior interval point of $A_\rho$ and a boundary half-disk point of each
adjacent sector. Similarly, $[g^\pm]$ is an endpoint of each $A_\rho$ and a
boundary point of each sector. These local claims use the actual rank-two
circle, not a homotopy description of a link.

We next prove that the other marked points on a meridian occur between its
two endpoints. Fix $[h]\in A_\rho$. Choose an old vertex $w\in h$, $w\neq v$.
Such a vertex exists: otherwise every old line would meet $h$ at $v$, so the
whole arrangement would be concurrent. Choose the signed cocircuit $Y_w$
with $(Y_w)_g=+$. The localization evaluation at $Y_w$, interpreted as
$-1,0,1$ and extended linearly over every order simplex, defines a continuous
map
\[
 e_w:A_\rho\longrightarrow[-1,1].
\]
A simplex cannot have both a positive and a negative value, since its
vertices are comparable localizations. The only vertex of $A_\rho$ at
which this evaluation vanishes is $[h]$: every extension in the fiber
already passes through $v$, and an ordinary pseudoline cannot also pass
through $w$ on the same old line $h$. The two parallel extensions with
underlying line $h$ have different images $\rho,-\rho$, so only one lies
in this meridian. Thus $e_w^{-1}(0)=\{[h]\}$, while its values at $[g^+]$
and $[g^-]$ have opposite signs. Each such marked point separates the two
endpoints.

Here is the elementary topological fact we use to cut a meridian. If a
contractible finite CW complex has a point with an open interval
neighborhood, cutting out a smaller interval gives exactly two contractible
components whenever the cut point separates two specified points. To prove
it, restore the interval as a one-cell joining the cut complex. Its relative
homology and the long exact sequence show that the cut complex has two
acyclic components. Van Kampen identifies the fundamental group of the
restored complex with the free product of the groups of those components,
so both are simply connected. A simply connected acyclic CW complex is
contractible, by Hurewicz followed by Whitehead. Iterating proves the
statement for finitely many disjoint cuts. Since every marked point in
$A_\rho$ separates its endpoints, their separating order is linear: cutting
one places every other point on one of the two endpoint sides. No point
can be on a branch outside the endpoint-to-endpoint order.

Choose disjoint sufficiently small polygonal disks about all $2n$ oriented
parallel points, compatible with the link sectors just described and
invariant under $\tau$. Let $P^\circ$ be the complement of their interiors,
and let $R_\gamma^\circ,A_\rho^\circ$ be the induced cut pieces. Each
$R_\gamma^\circ$ is contractible. Indeed, in $R_\gamma$ each removed piece
is a boundary half-disk (a sector at $g^\pm$ is homeomorphic to a half-disk),
and the half-disk strongly deformation retracts to its frontier semicircle
while fixing that frontier. Thus deleting it leaves a deformation retract
of the original sector. The meridian $A_\rho^\circ$ is a disjoint union of
contractible components: endpoint cuts do not change its homotopy type,
and the intervening interval cuts were handled in the previous paragraph.
The cuts at $g^\pm$ eliminate all intersections of nonadjacent sectors and
all triple intersections. All pieces and intersections can be triangulated
as finite subpolyhedra, so their inclusions are cofibrations.

We now construct the comparison map with its boundary marking. Start with
the ordinary suspension sphere $S^2=\operatorname{susp}|C|$, whose meridians
are indexed by $\rho$ and whose sectors are indexed by $\gamma$. On the
$\rho$-meridian place one marked point for each $[h]\in A_\rho$, in the
separating order just proved, and mark the two suspension tips. Remove small
disks about these $2n$ marked points. We obtain a sphere $S^\circ$ with
$2n$ boundary circles. The cut sectors are disks and the cut meridians are
intervals in exactly the same incidence diagram as
$R_\gamma^\circ,A_\rho^\circ$.

Map each boundary circle homeomorphically to the corresponding actual
circle in $P^\circ$, matching its two half-circle pieces at a non-tip
parallel point and its sector pieces at a tip. For each cut meridian
interval, its two endpoints have now been prescribed in the corresponding
contractible component of $A_\rho^\circ$; choose a path between them in
that component. These paths give compatible maps on all sector boundaries.
Each such boundary map extends over the corresponding sector disk because
$R_\gamma^\circ$ is contractible. Gluing gives a continuous map
$S^\circ\to P^\circ$ which is a homeomorphism on every marked boundary
circle. Its restriction to every sector is a homotopy equivalence between
contractible spaces, and its restriction to each intersection component is
likewise a homotopy equivalence. Gluing finite CW homotopy equivalences
along these cofibrations proves that the whole map is a homotopy
equivalence. One may verify this gluing assertion inductively by writing
each addition of a sector as a pushout along the disjoint union of its
intersection components; such a pushout is a homotopy pushout because all
inclusions are cofibrations.

All choices can be made equivariantly. The sign involution exchanges the
tips, sends $\rho$ to its antipodal vertex, and reverses each meridian's
endpoint order. On the model sphere choose the action $(z,t)\mapsto(-z,-t)$;
place paired marked points at opposite heights on antipodal meridians.
Choose the boundary maps, meridian paths, and sector fillings for one
representative of each antipodal pair and transport them by $\tau$. No
sector or meridian is invariant under this action, so these prescriptions
have no fixed-piece compatibility obstruction. The resulting equivariant
map descends to a homotopy equivalence of quotient pairs: its lift is a
homotopy equivalence, its induced map on the double-cover deck groups is
the identity, and the covering-space exact sequence and Whitehead give
the quotient conclusion. The model quotient is exactly
$\mathbb {RP}^2$ with $n$ open disks removed, not an unspecified quotient
of a homotopy sphere.

Finally attach a M\"obius band to every boundary circle on both sides of
this pair equivalence, using the same parametrizations. The boundary maps
are homeomorphisms and the attachments are cofibrations; the same homotopy
pushout argument extends the equivalence. On the model surface this is the
real point blow-up at each of $n$ points, giving the connected sum of
$n+1$ projective planes. Each replacement lowers Euler characteristic by
one, which also proves the displayed Euler formula directly.
\end{proof}

\subsection{The disk gluing, its local links, and its cell counts}

We use the later surface $\mathcal X_2$ of cote 154/005 and /007.
Its faces are the ordinary relative positions, and its additional edges
are the sweeps across an old pseudoline. This convention is the one used
by the count on 154/120. The capped version $\mathcal X_0$ and the versions
with a supercritical face are distinguished below. All our arrangements
are finite, labelled, and of rank three. An old intersection may have
arbitrary multiplicity. The boundary of the disk obtained by cutting
along an added pseudoline is its two-sheeted normal-orientation cover.
Its generization packets and empty specialization fans give the
transition intervals of 154/035--038. We shorten exceptional intervals,
as on /038, so distinct closed transition intervals are disjoint.

We call a complementary interval an \emph{inertia interval}, or a corner
after it is contracted in the polygonal cell model. Ordinary gluing uses
the reflection selected on 154/039: the two extreme chord names are
exchanged. The same-chord endpoint correspondence rejected on that page
is a different gluing and is not used here. For a multiple vertex, the
whole fan is treated at once; the two extreme strands determine the two
endpoints, and the internal strands determine its weight.

\begin{resultenv}{T1}{Proved}{The lens lemma and the only repeated-vertex corner}
Let $D$ be ordinary. If two consecutive generization intervals on its
disk boundary correspond to the same old vertex, then $D$ passes through
exactly one old vertex. The intervening side either contains an ordinary
specialization interval or bounds an exceptional corridor beside an old
pseudoline. Consequently the complete transition cycle, including the
exceptional sweep intervals of $\mathcal X_2$, never has a corner bounded
by two generizations at the same vertex.
\end{resultenv}
\begin{proof}
Different vertices on $D$ use disjoint old-line sets. In the disk, the
chords in one such set are disjoint and form two opposite packets. Every
chord belonging to another such set crosses each of them, so that other
set has one packet on each of the two complementary boundary intervals.
Thus the two packets of one vertex cannot be consecutive if there is a
second vertex on $D$.

Suppose now that $D$ passes only through $v$. Let $C$ be the parallel
class of chords coming from the old lines through $v$. On either side,
let $P$ be the boundary arc between the two $C$-packets, and take the
outermost chord $g\in C$ facing $P$. The arc $P$ and $g$ bound a disk
$U$. Every other old chord has exactly one endpoint on $P$ and crosses
$g$ exactly once. Its part in $U$ is an arc from $P$ to $g$. No other
chord of $C$ lies in the interior of $U$.

Orient these arcs from $P$ to $g$. Their intersection graph has no
directed cycle. Here is the precise known result giving this assertion:
decone the rank-three arrangement $\Sigma\cup\{D\}$ at $g$ and use
$D$ as objective. The chosen lune lies in one affine sign region
$D>0$, after reversing its sign if necessary. On this sector $U$, the directed edges
of its oriented-matroid program are exactly the indicated orientations
(reverse the objective sign if necessary). Every strand here is
nonparallel to the objective: a strand through $v$ would belong to $C$.
The graph therefore has no undirected edges in $U$.
Sturmfels and Ziegler's definition of this graph on pp.~30--31,
Proposition~3.9 on p.~33, and the rank-three Euclidean result,
Corollary~4.5 on p.~40, imply the asserted absence of directed cycles
\cite{SZ1993}. Restricting an acyclic directed graph to
the vertices and edges in $U$ preserves acyclicity.

If $U$ contains an old vertex, choose a source in this finite directed
graph. Along each incident strand, its incoming segment runs directly
from $P$ to that vertex: an earlier vertex would be a predecessor.
These incoming endpoints form a consecutive block on $P$. Indeed, a
strand beginning between the two extreme incoming endpoints and not
ending at the source must leave the intervening fan to reach $g$.
The Jordan separation theorem forces it to cross an incoming extreme
segment before the source, again giving a predecessor. The same
argument excludes a vertex inside the fan. Including every strand
through the source therefore gives an empty specialization fan with
base on $P$, hence an ordinary specialization interval between the
two generization packets.

If there is no such specialization interval, $U$ has no old vertex.
Its intervening strands are mutually disjoint, and the region between
$D$ and the old line underlying $g$ is exactly an empty corridor.
Conversely an empty corridor contains no old vertex that could be a
specialization apex on this side. The shortened exceptional transition
interval occupies its boundary side and separates the two generization
intervals. This proves the final assertion, including multiple vertices.
\end{proof}

\begin{resultenv}{T2}{Proved}{Ordinary corner circuits are four-cycles}
For the reflection gluing above, every vertex circuit using only
ordinary transition edges has four corners from four distinct relative
positions. It contains one gen--gen corner, one spec--spec corner, and
two mixed corners. This remains true when a position passes through
arbitrarily many old vertices or a semi-stable position has arbitrarily
many critical specializations.
\end{resultenv}
\begin{proof}
We first check a gen--gen corner. By T1 its two intervals belong to
distinct vertices $v,w$. Their old-line sets are disjoint. Choose
disjoint closed coordinate disks at $v,w$, containing no other old
vertex, in which the old pseudolines are radial strands and $D$ is a
diameter. The choice of one of the two lifted packets at each vertex
specifies one of its two resolutions. Denote the corresponding signs
by $\epsilon,\eta$. The four resulting extensions have the following
two localization coordinates; all their other coordinates agree:
\[
\begin{array}{c|c|c}
\text{position}&(\sigma(v),\sigma(w))&\text{corner type}\\\hline
D&(0,0)&G,G\\
D_v&(\epsilon,0)&S,G\\
D_{vw}&(\epsilon,\eta)&S,S\\
D_w&(0,\eta)&G,S.
\end{array}
\]
The resolutions commute because their supports are disjoint. Their
localizations differ, so the four labelled positions are distinct,
also modulo reorientation of the new pseudoline: a nonzero coordinate
outside the indicated zero set fixes that reorientation.

We verify that this square is the actual corner circuit, rather than
merely a commuting diagram of extensions. In a resolved fan the selected
generization packet becomes the base of the reverse specialization
fan. Its extreme strand names reverse order. The chosen reflection
exchanges those names, and thus identifies the same two \emph{spatial
sides} of the interval before and after resolution. On the selected
side each extreme incoming strand now encounters the old apex first.
It cannot acquire another specialization apex before it; every such
triangle would have to cross that first vertex. The portions of all
strands outside the resolving disk are unchanged. Hence the adjacent
inertia interval remains adjacent to the reverse transition on that
side. All possible changes at the other lifted packet lie outside
this inertia interval: before resolution that other packet was itself
a transition interval, and our two intervals were consecutive. This
argument applies to a fan of any multiplicity, using its two extremes.

An exceptional corridor cannot newly intrude into the tracked interval.
If one resolution produces a position near an old line through the
remaining zero vertex, the just resolved, different vertex lies in the
tracked lens and blocks emptiness on that side. Its empty corridor, if
present, is on the other side. If the old line does not contain the
remaining zero vertex, nearness to it is impossible. After both
resolutions, an exceptional wall in the tracked spec--spec gap would
have to contain both distinct apices $v,w$: these are the two fans
bounding that gap. No old line contains both, by the disjoint-support
argument. Thus no exceptional interval enters that gap either.
We may therefore perform the two resolutions in either order while
following the same corner. The local circuit is exactly the displayed
four-cycle. Its inverse endpoint maps return the starting germ to
itself, so no additional identification occurs after four steps.

It remains to prove that no other all-ordinary circuit has been missed.
The endpoint identifications form a finite permutation of corner germs,
so such a circuit is finite. Across an ordinary edge the integer
$\nu(D)$ changes by exactly one. Choose a position of maximal $\nu$
on the circuit. Both adjacent steps decrease $\nu$, so its corner is
gen--gen. T1 excludes a repeated-vertex gen--gen corner, and the preceding
local computation identifies its entire orbit with a four-cycle.
Thus every ordinary circuit has the stated form. Neither the argument
nor its conclusion bounds the total number of specializations at any
one position.
\end{proof}

\begin{resultenv}{T3}{Proved}{Exceptional completion, embeddings, and the surface counts}
The prescribed $\mathcal X_2$ is a well-defined compact connected
nonorientable surface, up to labelled cellular homeomorphism. Every
vertex has valence four, and every closed face, edge, and vertex embeds.
With $b=\sum_i n_i$ and $F(t)=\sum_j c_jt^j$ as above, its cell counts are
\[
 s_2=F(1),\qquad s_1=2F'(1)+2b,\qquad s_0=F'(1)+b.
\]
Its ordinary vertices are the gen--gen inertia intervals; its exceptional
vertices are the $2b$ flags $(v,e)$ of the original arrangement.
\end{resultenv}
\begin{proof}
We make the exceptional construction explicit. Fix an old pseudoline
$D_i$, with cyclically ordered vertices and edges. In a sufficiently
small Moebius neighborhood of $D_i$, an ordinary pseudoline close to
$D_i$ crosses it once. Its crossing is either in an old edge $e$, or at
an old vertex $v$. In each case there are two choices of the local
normal direction. We denote these positions $U_e^+,U_e^-$ and
$S_v^+,S_v^-$. The superscripts are local normal labels; transporting
them once around $D_i$ exchanges them. There are exactly $2n_i$ stable
positions $U_e^\pm$ and $2n_i$ semi-stable positions $S_v^\pm$ near
$D_i$. To verify exhaustiveness, cut the neighborhood at its single
crossing with $D_i$. On the resulting strip, a close pseudoline is a
section with fixed side until that crossing and the opposite side
after it. Contracting each old-vertex fan leaves no further choice.
Conversely these sections realize all the indicated positions.

If a position admits an empty corridor with $D_i$, it is one of these
positions. Slide its boundary through the corridor; there is no old
vertex in the interior, so no localization changes before reaching
a sufficiently small neighborhood. This is an explicit isotopy
relative to the old arrangement's combinatorial cells, and gives
the required comparison with the near positions. It introduces no
additional relative-position class.

The exceptional edges pair $U_e^+$ with $U_e^-$, and $S_v^+$ with
$S_v^-$. There is one such sweep for every unoriented edge and vertex
of $D_i$, hence $2n_i$ edges. At every original flag $v\in\partial e$
the four incident positions, in cyclic order, are
\[
 U_e^+,\quad S_v^+,\quad S_v^-,\quad U_e^-.
\]
The first and third sides are ordinary specialization/generization
sides, and the other two are exceptional sweeps. The endpoint rule
preserves the original flag $(v,e)$. On a stable sweep it preserves
the old vertex at that end; on a semi-stable sweep it preserves the
old edge on that side. These are the two equivalent descriptions of
the same strip gluing. The four positions are distinct: the stable
ones have no zero, the semi-stable ones have exactly $v$ as zero,
and the two normal choices differ on the neighboring fan signs.
A vertex off $D_i$ exists because the arrangement has rank three;
its fixed nonzero sign prevents these differences from being removed
by reorienting the new element. Thus each exceptional circuit is a
four-cycle with four distinct face labels.

There is no missing corner between two exceptional intervals. At a
stable near position, an exceptional interval is bounded by the two
ordinary fans at the distinct endpoints of its old edge. Two different
old wall lines cannot occupy that same gap: both would have to contain
those two vertices. At a semi-stable near position, the exceptional
interval lies between the two packets of its single vertex, and the
outermost wall of the lens on that side is unique. Thus every corner
meeting an exceptional edge is one of the flagged circuits just listed.

We also check exhaustiveness of the boundary circuits before inserting
any exceptional transition intervals or sweeps. The following maximum
argument is on those ordinary inertia circuits. A circuit not capped by T2 has a maximum of
$\nu$; its gen--gen corner must therefore use the same vertex, and
T1 identifies a specific empty corridor with an old line $D_i$.
Resolving its vertex toward an incident old edge $e$ gives $U_e$.
The corridor remains empty, and the next reverse fan is at the other
endpoint of $e$, giving $S_w$. Following the boundary thus moves along
the vertices and edges of the same $D_i$, without changing its label.
Conversely every such old flag supplies this local step. After one
turn the normal sign is reversed, because a pseudoline in the
projective plane is one-sided; after two turns it returns. This proves
that the boundary circuits are exactly the normal-orientation circles
of the old lines, with $4n_i$ successive stable/semi-stable sectors.

One can check the twisting without any orientation inference from the
counts. Orient all near positions as generizations of $[i^+]$, using
a nonzero cocircuit off $D_i$ to make this unique. Lemma~4.1 of
Sturmfels--Ziegler identifies their cyclic incidence with
$\mathcal E(M/i)$. The sweep sends
$\kappa_+(\rho)$ to $\kappa_+(-\rho)$: it reverses just the signs at
the vertices of $D_i$ and preserves the signs off $D_i$.
Thus the actual pairing of the free boundary intervals is the
antipodal identification on this circle. The quotient of its collar
by $(z,0)\sim(-z,0)$ is a Moebius band.

In explicit coordinates the inserted band is
$([0,1]\times[-1,1])/((0,t)\sim(1,-t))$.
Its boundary is that normal-orientation double cover. Put fiber edges
in the dual cell picture over the vertices and edge-midpoints of the
old $D_i$. The intervening
rectangles are indexed by the old flags and have exactly the four
corner labels displayed above. Thus the prescribed exceptional
endpoint pairing is the actual Moebius-band gluing, rather than a
conclusion drawn from an Euler-characteristic calculation.
Filling each such circuit by a disk gives the capped surface
$\mathcal X_0$. Replacing that disk by a Moebius band, with the flag
subdivision just described, gives exactly $\mathcal X_2$; this is the
topological operation stated on 154/007. In particular every boundary
component is filled, and the inserted Moebius band is an embedded
subsurface. This proves nonorientability directly, without relying on
the manuscript's monodromy assertion on 154/064.

For completeness, we check the topology and the embedding issue.
First glue the original disks only along their pairwise disjoint
closed transition intervals. Each interior point of such an interval
has exactly two representatives, and no interval endpoint belongs
to a second transition interval in the same disk. Thus no zigzag
of these identifications joins distinct points of one original disk.
The finite no-zigzag criterion SA10--A4.3, rechecked here, gives
embeddings. In the polygonal cell model, equivalence classes of corner
germs are exactly the four-cycles of T2 and the preceding exceptional
calculation. Each has four distinct disk labels. Hence no two corners
of a single normalized face are identified, no edge has its endpoints
identified, and no edge interior is identified with any other part
of the same face. The same no-zigzag criterion proves that all closed
cell characteristic maps embed. One may equivalently retain the
original disks and distribute the four sectors of each filled inertia
circuit among its incident faces; this gives the same closed-cell
model without collapsing an interval in the original disk.

The quotient relation on the finite union of closed polygonal disks is
closed: it is the finite union of paired edge graphs, diagonal interior
relations, and finite corner classes. The quotient is compact Hausdorff.
A face-interior point has a disk neighborhood, an edge-interior point
has two half-disks, and a vertex has the cone on a genuine four-cycle,
hence a disk neighborhood. This proves the surface assertion by actual
local links, as required by the distinction in SA10--C3.04.

The surface is connected. The cubical part of $T(M)$ obtained by
removing the interiors of the parallel polygonal cells is connected:
the full complex is connected by E1, and attaching a two-cell along
its connected boundary cannot join two components. Its face-incidence
graph is exactly the graph of ordinary positions and ordinary
transitions. That graph is therefore connected, and adding the
exceptional sweeps preserves connectedness.

There is one face per ordinary position, giving $s_2=F(1)$. A position
with $j$ zero vertices has exactly $2j$ immediate generizations,
each counted once at its more special endpoint. Thus there are
$2F'(1)$ ordinary edges. The exceptional calculation adds $2b$ edges.
Every vertex is four-valent, so $2s_1=4s_0$, giving the displayed
counts. Each ordinary four-cycle has exactly one gen--gen corner by
T2, and each exceptional vertex has exactly one original flag, proving
the final indexing assertion.

Finally, the construction depends only on the labelled chord diagrams
and the specified endpoint permutations. Changes of interval lengths
or parametrizations are absorbed by boundary homeomorphisms, which
extend across each disk by coning. Choice of a starting point or
orientation of a cyclic word merely reparametrizes the same polygon.
If the labelled internal chord graph is retained, match its vertices
and edges first, and extend across each complementary disk cell;
this gives the same extension while preserving that graph. These
extensions give a labelled cellular homeomorphism of the quotients.
This is the precise ambiguity in ``well-defined.''
\end{proof}

\paragraph{Scope of the edge cases.}
A rank-three arrangement of distinct pseudolines has $n\geq3$ and at
least two vertices on every old line. Thus neither an empty arrangement,
a one-line arrangement, nor a pencil satisfies the hypotheses. Multiple
intersections are included throughout T1--T3. All cells here are connected
disks or intervals; disconnected strata, infinite arrangements, and
non-locally-finite gluing families are not asserted to satisfy these
finite statements. No faithfulness or divisibility axiom from the
separate tame-space questions enters the construction.
\subsection{The topology formula and the refutation of L1}
The hypotheses in E5 below all hold for the surface of this paper, by T3.
\begin{resultenv}{E5}{Proved under stated hypotheses}{The numerical correction forced by the cubical extension counts}
Suppose the manuscript surface has the verified cell counts
\[
 s_2=F(1),\qquad s_0=F'(1)+\sum_i n_i,
\]
and is a closed connected nonorientable surface. Then
\[
 \chi(\mathcal X)
 =1-n-4\sum_{k\geq1}k\,c_{2k+1},
 \qquad
 \mathcal X\cong\#^{\,n+1+4\sum_{k\geq1}k c_{2k+1}}\mathbb {RP}^2.
\]
It is homotopy equivalent to $B(M)$ if and only if no ordinary relative
position passes through three old vertices. This hypothesis is automatic
for $n\leq5$.
\end{resultenv}
\begin{proof}
All the stated surface hypotheses hold for $\mathcal X=\mathcal X_2$ by
T3. They are not consequences of the extension-space theorem alone.
Under them $\chi(\mathcal X)=F(1)-F'(1)-F'(-1)$ by E3. Subtract E2 and
compare coefficients. For the coefficient of $c_j$ the difference is
\[
 1-j-j(-1)^{j-1}-(-1)^j=(j-1)((-1)^j-1).
\]
This is zero for even $j$ and $-2(j-1)$ for odd $j$, proving the formula.
The classification of compact surfaces gives the homeomorphism assertion.
By E4, $B(M)$ has the homotopy type of the connected sum of $n+1$ projective
planes. These two surfaces are homotopy equivalent exactly when their Euler
characteristics agree, which here means $c_{2k+1}=0$ for every $k\geq1$.
If any ordinary position has $\nu\geq3$, resolving all but three of its
vertices produces an ordinary position with exactly three zeros, so $c_3>0$.
Conversely $c_3=0$ excludes all such positions. Finally distinct vertices on
an ordinary added pseudoline involve disjoint old-line classes, each of size
at least two; hence $2\nu\leq n$ and $\nu\geq3$ requires $n\geq6$.
\end{proof}

\subsection{Certified enumeration and the six-pseudoline test}
\label{sec:counts}

\paragraph{Statement C.01 --- exhaustive encoding for a simple arrangement. (Proved)}\label{C.01}
Let $M$ be a simple, uniform oriented matroid of rank three on
$[n]=\{1,\ldots,n\}$, with a fixed chirotope $\chi_M$.
For a non-loop, non-parallel, same-rank extension by an element $p$, put
\[
 q_{ij}=\chi(i,j,p)\quad (1\leq i<j\leq n).
\]
The zero pairs of $q$ form a matching in the complete graph on $[n]$.
Conversely, we obtain every such extension by enumerating matchings,
assigning signs to the remaining pairs, and retaining precisely the
assignments satisfying the chirotope axioms on $[n]\cup\{p\}$.
We quotient by reorientation of $p$ by requiring the first nonzero
coordinate of $q$, in lexicographic pair order, to be positive.

\emph{Proof.}
If $q_{ij}=q_{ik}=0$ for distinct $i,j,k$, then $p$ lies in both
rank-two flats spanned by $\{i,j\}$ and $\{i,k\}$. Their union has
rank three because $M$ is uniform. Submodularity bounds the rank of
their intersection by one. That intersection contains $i$ and $p$,
so the non-loop $p$ is parallel to $i$, contrary to the hypothesis.
Thus the zero pairs are disjoint. Every extension therefore occurs in
our candidate list. We define the remaining triples by alternation,
set repeated triples to zero, and leave triples in $[n]$ equal to
$\chi_M$. The full signed basis-exchange axiom is then a necessary
and sufficient chirotope test, so accepting exactly its solutions
introduces no extra candidates. Multiplication of all $q_{ij}$ by
$-1$ is exactly reorientation of $p$. No allowed vector is zero, and
hence each orbit has exactly one representative with the stated
normalization. \hfill$\square$

For reproducibility, our fast enumerator first checks the necessary
three-term Grassmann--Pl\"ucker sign relations
\[
 \bigl\{
 \chi(a,b,c)\chi(a,d,e),
 -\chi(a,b,d)\chi(a,c,e),
 \chi(a,b,e)\chi(a,c,d)
 \bigr\}.
\]
Each such set must either be zero or contain both signs. This is
\emph{only a filter}: a second, independent program checks the full
signed basis-exchange axiom on every surviving assignment. Explicitly,
for every ordered basis $x=(x_1,x_2,x_3)$ and every ordered basis
$y=(y_1,y_2,y_3)$ it verifies that at least one $i$ satisfies
\[
 \chi(y_i,x_2,x_3)
 \chi(y_1,\ldots,y_{i-1},x_1,y_{i+1},\ldots,y_3)
 =\chi(x_1,x_2,x_3)\chi(y_1,y_2,y_3).
\]
It suffices to choose $x_1$ in each unordered basis, put the other two
entries in increasing order, and use increasing $y$: permutations of
the remaining entries multiply both sides by the same sign. We thus
check the full axiom, not merely its three-term consequence. For
$n=6$, the number of candidates, before this filtering, is
\[
 \sum_{k=0}^{3}
 \frac{6!}{(6-2k)!\,2^k k!}\,2^{15-k-1}=354304.
\]
The programs are \texttt{counts\_extensions.py}, its C++ inner loop
\texttt{counts\_extensions.cpp}, and
\texttt{counts\_independent\_check.py}. Their input and complete accepted
sign-vector lists accompany the paper.

\paragraph{Statement C.02 --- explicit input satisfying the test-case hypotheses. (Proved)}\label{C.02}
In homogeneous coordinates $[x:y:z]$ on $\mathbb{RP}^2$, take the six
lines with coefficient vectors
\[
\begin{split}
 a_1&=(1,0,1),\quad a_2=(2,0,1),\quad
 a_3=(0,1,1),\\
 a_4&=(0,3,1),\quad a_5=(1,1,1),\quad
 a_6=(4,4,1).
\end{split}
\]
Their twenty determinants, in lexicographic triple order, are
\[
\begin{split}
(&1,3,1,4,-2,-1,-7,-1,-13,-3,\\
 &-4,-1,-10,1,-14,-6,-2,-8,3,9).
\end{split}
\]
They are all nonzero; the arrangement is therefore simple and
nontrivial. The line $D=\{z=0\}$ contains
\[
 v_{12}=[0:1:0],\qquad
 v_{34}=[1:0:0],\qquad
 v_{56}=[1:-1:0].
\]
These are three distinct vertices. For $p=(0,0,1)$, its extension
vector in lexicographic pair order is
\[
 (0,+,+,+,+,+,+,+,+,0,-,-,-,-,0).
\]
All these claims follow by the displayed integer determinants or
cross products, which the attached programs also check.

\paragraph{Statement C.03 --- an independent disk-boundary certificate. (Proved under stated hypotheses)}\label{C.03}
We count the literal gen--gen inertia arcs of cote 154/120 independently
of the formula $s_1=2s_0$. Here are the complete rules used by the
checker, including the endpoint convention.

Realize the rank-two chirotope $q$ by oriented rays $w_i$ in an
oriented plane and order the $2n$ signed rays $\epsilon w_i$
cyclically. Equal signed rays form packets. For a two-element packet
$(\epsilon,i),(\eta,j)$, choose $k$ outside its unoriented parallel
class and order the first label before the second exactly when
\[
 -\eta\,\chi_M(i,j,k)q_{ik}=+1.                 \tag{C.1}
\]
The rank-two sign order and (C.1) suffice to construct the complete
boundary word. The global reversal of this word has no effect on
any count or gluing.

To verify (C.1), use local affine representatives
$a_i=(w_i,z_i)$ and $a_j=(\lambda w_i,z_j)$. At the endpoint in
direction $J(\epsilon w_i)$ of a sufficiently large circle,
the angular displacement has sign and relative order determined by
$\epsilon z_i/\lVert w_i\rVert$. Moreover
\[
 \det(a_i,a_j,a_k)
 =-(z_j-\lambda z_i)\det(w_i,w_k).
\]
For equal signed rays $\operatorname{sgn}\lambda=\epsilon\eta$;
substitution gives (C.1). This calculation is local to the
restriction on $\{i,j,k,p\}$ and does not assert a realization of
the whole extension. The two packets over the same vertex have
reversed orders, as required by the disk construction.

Each zero pair gives two generization intervals $G$, one in each
packet. For two adjacent endpoint labels $(\epsilon,i),(\eta,j)$
that belong to no $G$ packet, the interval is an ordinary
specialization interval $S$ precisely when, for every
$k\notin\{i,j\}$,
\[
 \chi_M(i,j,k)q_{ij}
     =\epsilon q_{ik}=\eta q_{jk}.             \tag{C.2}
\]
Indeed, the first expression is the side of line $k$ containing
$v_{ij}$ in the affine chart $p>0$. The other expressions are its
sides at the two endpoint rays. Agreement says that $k$ cuts neither
segment from $v_{ij}$ to those endpoints. Since two pseudolines
cross at most once and the boundary endpoints are adjacent, this
is equivalent to the intervening triangle being empty. Its
collapse is exactly the permitted specialization. The independent
checker confirms that the resulting pairs are exactly those
obtained by replacing one nonzero coordinate of $q$ by zero in
the exhaustive extension list.

For a possible exceptional sweep at $D_i$, the criterion is
\[
 q_{ab}\chi_M(a,b,i)=r\in\{+1,-1\}
 \quad\hbox{for every }a,b\ne i.              \tag{C.3}
\]
The common sign says that every base vertex off $D_i$ lies on the
same side of $D_i$ in the affine chart $p>0$. The other side is an
empty corridor. Its boundary interval lies on the side with sign
$-r$, and is shortened past the adjacent ordinary transition
intervals. Its position among the intervals is obtained from that
boundary subarc; the program uses its midpoint. When $q$ has one
zero pair, (C.3) is tested at the two lines through that vertex and
inserts the exceptional interval between the appropriate two
$G$ intervals. When $q$ has at least two zero pairs, there is no
such exceptional interval, since a zero pair avoiding $i$ would
violate (C.3).

These rules are exhaustive: the manuscript's ordinary transitions
are exactly the empty triangles and the parallel packets
(154/035--038), and its exceptional transitions are exactly the
empty corridors (154/036--038, together with the stable sweeps in
$\mathcal X_2$, 154/005). Consecutive $G$ intervals in the complete
cyclic list are therefore precisely the gen--gen inertia arcs.
Every boundary word and transition interval is recorded in the
boundary-certificate JSON files.

\paragraph{Statement C.04 --- the finite surface certificates. (Proved)}\label{C.04}
For the arrangement of C.02 the numbers $c_\nu$ of non-supercritical
relative positions through exactly $\nu$ vertices are
\[
 (c_0,c_1,c_2,c_3)=(80,150,75,10).
\]
In particular $s_2=315$. The direct boundary count of C.03 gives
$240$ gen--gen arcs at positions with two vertices and $60$ at
positions with three vertices. There are none at stable or
semi-stable positions. The $60$ flags of type $(0,1)$ in the
original simple arrangement contribute the exceptional vertices.
Thus the literal count in 154/120 is
\[
             s_0=240+60+60=360.               \tag{C.4}
\]

We also independently construct the complete edge identifications,
rather than infer the links from this count. An ordinary edge joins
an $S$ interval to the $G$ interval of its specialization with the
same unordered signed endpoint labels. The endpoint map exchanges
the two chord names. This is the reflection convention selected on
154/039, where one chord continues into the other chord; it is the
alternative to the discarded same-chord continuation. An exceptional
sweep at $i$ replaces $q$ by the normalized sign vector
\[
 \operatorname{sw}_i(q)_{ab}=
 \begin{cases}
   -q_{ab},&i\in\{a,b\},\\
   q_{ab},&i\notin\{a,b\}.
 \end{cases}                                  \tag{C.5}
\]
This is the local passage from one side of $D_i$ to the other.
For a stable exceptional edge, its endpoint map preserves the base
vertex at each end of the associated original edge. The two
adjacent $S$ intervals identify those vertices. For an unstable
exceptional edge, its endpoint map preserves the adjacent original
edge: resolve the adjacent $G$ interval to a stable position and
match its sweep partner on the other side. These are exactly the
four positions at an original vertex--edge flag.

The program \texttt{counts\_gluing.py} records every polygon, edge
pairing and endpoint permutation. A polygon corner is a consecutive
pair of transition arcs. Taking equivalence classes under the edge
endpoint maps yields $360$ classes. Every class contains exactly
four corners from four distinct polygons, and its link graph has
four distinct edges and degree two at each of its four vertices.
Thus every vertex link is a genuine four-cycle. Each edge is
paired once with an edge in a different polygon. No polygon has
two corners in the same equivalence class. Consequently each
closed polygon, edge and vertex embeds in the quotient; no
unverified closed-colimit embedding assertion is used. The finite
gluing relation is closed, so the quotient is compact Hausdorff.
The interior, edge, and vertex neighborhood tests prove that it is
a closed surface.

The adjacency graph of the polygons is connected. The signed
orientation constraints on adjacent polygons are inconsistent;
the full certificate contains an explicit orientation-reversing
cycle. These facts are checked directly from the endpoint maps.
There are $660$ ordinary edges and $60$ exceptional edges, hence
\[
 (s_0,s_1,s_2)=(360,720,315),\qquad
 \chi(\mathcal X)=-45.
\]
By the classification of compact connected surfaces,
$\mathcal X\cong\#^{47}\mathbb{RP}^2$ for this arrangement.
In particular the claim $\chi(\mathcal X)=1-n$ is false.

For clarity about the finite certification, every candidate is
specified by an explicit sign vector, every accepted candidate
satisfies the full chirotope axiom, and every gluing is an explicit
finite permutation of edge endpoints. The asserted counts and
link properties are exhaustive computations on those lists, with
assertions verifying each condition. The full axiom checker makes
$1\,089\,945$ signed basis-exchange checks for C.02. The computation
is therefore a complete finite counterexample, not a statistical
search or an inference from a partial enumeration.

\paragraph{Statement C.05 --- dependence on the arrangement at fixed $n$. (Proved)}\label{C.05}
As a second six-line input, take coefficient vectors
$(t,t^2,1)$ for $t=0,1,\ldots,5$. The same complete procedure gives
\[
 (c_0,c_1,c_2,c_3)=(74,132,57,4),\qquad
 (s_0,s_1,s_2)=(288,576,267).
\]
All $288$ vertex links are four-cycles, all closed faces embed,
and the quotient is connected and non-orientable. Thus this
surface is $\#^{23}\mathbb{RP}^2$, with Euler characteristic $-21$.
The independent axiom check makes $928\,233$ signed basis-exchange
checks. Together with C.04, this proves that the topology is not
a function of $n$ alone.

\begin{center}
\begin{tabular}{lrrrrrr}
\hline
Input & $c_0$ & $c_1$ & $c_2$ & $c_3$ & $s_0$ & $s_2$\\
\hline
Three lines & 4 & 6 & 0 & 0 & 12 & 10\\
Four convex dual points & 12 & 18 & 3 & 0 & 36 & 33\\
Five convex dual points & 31 & 50 & 15 & 0 & 100 & 96\\
Six convex dual points & 74 & 132 & 57 & 4 & 288 & 267\\
Six-line input C.02 & 80 & 150 & 75 & 10 & 360 & 315\\
\hline
\end{tabular}
\end{center}
The first three rows reproduce the earlier small cases and extend
the test to one five-line input. They are not substituted for a
claim about arbitrary arrangements.

\paragraph{Statement C.06 --- the fixed straight-line subproblem differs. (Proved)}\label{C.06}
For the particular realization C.02, the fifteen dual lines have
$42$ ordinary double vertices, one ordinary triple vertex at
$p=(0,0,1)$, and six supercritical vertices of multiplicity five.
Exact cross products give
\[
 V=49,\qquad E=42\cdot2+3+6\cdot5=117,\qquad
 F=1-V+E=69.
\]
The non-supercritical \emph{straight} positions therefore number
$(V-6)+E+F=229$, and their sector count is $2E=234$.
Their corresponding surface has Euler characteristic $-5$.
These counts are independently produced by
\texttt{counts\_straight\_dual.py} using integer arithmetic.
They do not equal the pseudoline counts $315$ and $360$.
Stretchability after moving the original lines cannot identify
all pseudoline extensions with straight extensions of a fixed
realization. In particular, around the displayed triple position
there are eight topological stable resolutions, while the fixed
straight realization supplies only six angular sectors.

\paragraph{Statement C.07 --- the first size allowing three-vertex positions. (Proved)}\label{C.07}
If a non-supercritical pseudoline $D$ passes through $\nu$ vertices
of a nontrivial arrangement of $n$ distinct pseudolines, then
$2\nu\le n$. Each such vertex uses at least two original lines;
no original line can occur at two different vertices on $D$,
because it would meet $D$ twice. The sets of original lines used
at those vertices are therefore disjoint. This proves the
inequality and shows that $n=6$ is the first possible size for
a three-vertex position. C.02 realizes this bound.
\section{A.2. The cellular comparison and Lead L2}
\begin{resultenv}{E6}{Refuted}{What the cyclic comparison actually selects}
Lead L2 is false: all $2^\nu$ stable generizations are present in
$\mathcal X$. Its true replacement is the following explicit cellular
description. Subdivide every ordinary cube of $\overline T(M)$ by its
coordinate midplanes. Its vertices are the centers of the original cube
faces, hence ordinary extension labels; its edges are the ordinary cover
relations. Replace each parallel polygon by the M\"obius band of E4,
subdivided as described below. In this CW model of $B(M)$ the dual
quadrangulation $\mathcal X^\vee$ is an embedded subcomplex containing
the entire one-skeleton.

Specifically, if $\sigma$ has distinct zero coordinates $i,j$ and resolving them
with specified signs gives $\sigma_i,\sigma_j,\sigma_{ij}$, its diamond is
\[
 \sigma<\sigma_i<\sigma_{ij},\qquad
 \sigma<\sigma_j<\sigma_{ij}.
\]
The associated quadrangle is the union of these two closed order-complex
triangles. It occurs in the ordinary part of the dual manuscript map
exactly when the specified transition ports at $i,j$ are adjacent in the
full boundary cycle of its disk. This construction specifies the ordinary
quadrangles and their attaching maps. The exceptional part consists of
the entire subdivided M\"obius bands. Their rectangular cells are indexed
by the original flags $(v,e)$, and their four vertices are
$U_e^+,S_v^+,S_v^-,U_e^-$ in that cyclic order.

Thus there are specified comparison maps
\[
 \mathcal X\cong\mathcal X^\vee
 \xrightarrow{\ \iota\ }B(M)
 \xrightarrow{\ \beta\ }\overline T(M).
\]
Here $\iota$ is a cellular subcomplex inclusion and induces a surjection on
fundamental groups; $\beta$ is the continuous blow-down collapsing each inserted M\"obius core circle
to its parallel point. The inclusion $\iota$ is a homotopy equivalence
exactly when no ordinary position passes through three old vertices.
In that case it is an identification of the indicated CW spaces.
\end{resultenv}
\begin{proof}
The independent-resolution argument of E2 proves that all $2^\nu$ stable
generizations exist. Cote 154, scan~45, explicitly indexes faces by all
relative positions except the supercritical ones; there is no stable-state
selection in this definition. An explicit position with $\nu=3$ therefore
has eight stable faces, not six. The six-line counterexample and its finite
certificate in Section~\ref{sec:counts} instantiate this discrepancy.

We first identify the subdivision precisely. Give the cube of an ordinary
localization $\sigma$ its coordinates in $[-1,1]^{\nu(\sigma)}$, one for
each zero cocircuit pair. A coordinate remains free on a face if its
zero is unresolved, and equals $+1$ or $-1$ if that zero has been
resolved. Subdivision at coordinate zero therefore gives vertices
indexed by the faces of this original cube. A small $k$-cube is indexed
by an interval of length $k$ in its face order: its least extension has
$k$ specified zeros that may be resolved, and its greatest extension
resolves all of them with the indicated signs. These coordinate
descriptions restrict identically on every common original face, so
the subdivisions glue over the whole ordinary cubical subcomplex.
In particular the edges are exactly cover relations, and the small
two-cubes are exactly incidence diamonds. Triangulating each diamond
along its least-to-greatest diagonal gives its two order simplices.

For the assertion about selection, the two displayed chains are distinct
simplices of the order complex. Their common edge is
$[\sigma,\sigma_{ij}]$, their other edges give the four sides of the
diamond, and their union is an embedded closed disk. In the disk-gluing
construction, a corner bordered by the two indicated transition intervals
is traversed by performing these two independent local resolutions in the
two possible orders. Their supports are disjoint, so the final extension
is the same $\sigma_{ij}$. Conversely an ordinary corner of this type
contains consecutive transition ports, by the definition of a corner of
the disk. This gives the criterion and the attaching map.

No interior simplex is counted twice: a triangle in such a diamond records
its unique least extension, the intermediate resolution, and the final
two-coordinate resolution. The other intermediate vertex is then uniquely
determined. Distinct diamonds therefore meet only in their common faces.
Their union, with their faces, is a simplicial subcomplex after adding the
indicated diagonals. The sign involution carries each diamond to its
opposite diamond, and its quotient remains embedded by the sign argument
in E2.

We now include the exceptional cells, so the description gives the whole
dual surface. Fix an old line $D_i$ and its corresponding parallel
polygon in $\overline T(M)$. Lemma~4.1 of Sturmfels--Ziegler identifies
its boundary with the rank-two circle $\mathcal E(M/i)$, with its order
reversed. After the subdivision just made, denote this boundary circle
by $C_i$. Its $4n_i$ vertices alternate between $S_v^\pm$ and $U_e^\pm$,
and its edges are the ordinary transitions at the flags $v\in\partial e$.
The rank-two sign involution $a_i:C_i\to C_i$ exchanges the two superscripts
and preserves this subdivision. It is a fixed-point-free half-turn.
The interval bundle
\[
 H_i=(C_i\times[-1,1])/\bigl((z,t)\sim(a_i z,-t)\bigr)
\]
is a M\"obius band. Its boundary is identified with $C_i$ by
$z\mapsto[z,1]$. The fiber over the pair $\{z,a_i z\}$ joins these two
boundary vertices. Taking one such fiber edge for every antipodal
vertex pair gives $2n_i$ edges. The complementary closed rectangles,
one per antipodal edge pair of $C_i$, give $2n_i$ two-cells. Each is
the rectangle with corners
$U_e^+,S_v^+,S_v^-,U_e^-$ for one flag $(v,e)$.

The corridor isotopy and endpoint calculation in T3 identify these
fiber edges with the exceptional sweep transitions and these rectangles
with their four-corner circuits. In particular this is equality of
the attaching maps on every boundary edge, not merely a numerical
comparison. Every $C_i$ embeds in the ordinary complex by E2; the band
interiors are disjoint, and each band rectangle is a literal embedded
rectangle in $H_i$. The nonconcurrency hypothesis gives $n_i\geq2$, so
the four vertices of each rectangle are distinct. Consequently the
ordinary selected diamonds together with all the $H_i$ form an embedded
CW subcomplex of the specified $B(M)$. T2 and T3 supply the complete
ordinary and exceptional corner circuits, respectively, so the cell
correspondence is bijective and respects every attaching map. It gives
the asserted cellular homeomorphism from $\mathcal X^\vee$ to that
subcomplex. This also proves that no exceptional cells have been omitted.

The inclusion contains the entire one-skeleton: every ordinary cover
relation occurs as a transition, and all band fiber edges occur as
exceptional transitions. Building $B(M)$ from this subcomplex requires
only the remaining ordinary cells of dimension at least two. Their
attachment induces a quotient on fundamental groups, proving that
$\iota_*$ is surjective. More explicitly, its kernel is the normal
subgroup generated by the boundary loops of the unselected incidence
diamonds: these are precisely the additional two-cells, and attaching
cells of dimension at least three does not change the fundamental group.
Paths from a fixed base point to the attaching loops change their
representatives only by conjugation, so this normal-subgroup description
has no path ambiguity. The additional $k$-cells for $k\geq3$ are all the
ordinary small cubes indexed by length-$k$ face intervals above. In total
the ordinary subdivision has
\[
 N_k=2^k\sum_{j\geq k}\binom jk c_j
     =\frac{2^k}{k!}F^{(k)}(1)
\]
$k$-cells: inside an original open $j$-cube choose its $k$ varying
coordinates and their $2^k$ half-intervals, keeping the other coordinates
zero. Therefore $B(M)$ is obtained from the embedded surface by attaching
$2F''(1)-F'(1)+b$ unselected squares and $N_k$ cells of each dimension
$k\geq3$, where $b=\sum_i n_i$. Here we used T3: the surface has
$F'(1)-b$ ordinary squares and $2b$ band squares. This is a finite
relative CW description of the specified comparison, including every
attaching map through its cube-face interval.

For the blow-down map choose an equivariant
homeomorphism $h_i:C_i\to S^1$, so that
$h_i(a_i z)=-h_i(z)$. In the characteristic disk of the parallel cell
define
\[
 \beta_i([z,t])=t\,h_i(z).
\]
This is well-defined, is the specified boundary identification at
$t=1$, and collapses precisely the core circle at $t=0$. Extend it by
the identity on the ordinary complex. The resulting map $\beta$ has
exactly the stated target and agrees with the ordinary comparison on
every shared edge. Equivalently, $\overline T(M)$ is the quotient of
$B(M)$ by its $n$ disjoint core circles. In particular $\beta_*$ is
surjective and its kernel is the normal subgroup generated by these
core loops, by van Kampen.

Suppose finally that no ordinary position has three zero vertices.
Every ordinary two-cube has exactly four generization ports, no
ordinary specialization port, and no exceptional port. The last claim
follows from T3, because a position near an old line has at most one
zero vertex. T1 makes the four ports alternate between the two zero
vertices, so all four of the subdivided square's diamonds are selected.
There are no higher-dimensional ordinary cubes. Thus the subcomplex
is the whole of $B(M)$, proving the identification in this case. If a
position has at least three zeros, E5 gives a strict Euler-characteristic
difference, so $\iota$ cannot be a homotopy equivalence.

There are $2\nu$ generization ports at a position with $\nu$ zeros. When
there are no intervening specialization ports, their cyclic link has
$2\nu$ adjacent pairs. In general the relevant cycle includes the
specialization ports too; the rule is adjacency in that full cycle. Neither
count is a count of the stable vertices of the $\nu$-cube, which number
$2^\nu$. Furthermore the quadrangles live in its face-incidence subdivision:
their vertices are cube-face centers of dimensions $\nu,\nu-1,\nu-1,\nu-2$.
They are not the two-dimensional faces of the original unsplit cube.
The unmodified extension quotient also does not contain the exceptional
bands by this comparison: its blow-down collapses each entire core
circle. Thus the full cellular embedding belongs to the specified real
blow-up $B(M)$, while its ordinary selected diamonds lie in
$\overline T(M)$ itself.
\end{proof}

\paragraph{Literature scope.}
Sturmfels--Ziegler Theorem~2.3, p.~28, identifies the extension poset of a
\emph{realizable} rank-$r$ oriented matroid with a Baues poset for an
$n$-cube projected to an $(n-r)$-dimensional zonotope. That theorem is an
isomorphism of posets; it neither defines the above cyclic port rule nor
asserts that the manuscript surface is a coherent subcomplex. The rank-four
disconnected examples of Mn\"ev and Richter-Gebert (1993, pp.~271--273)
\cite{MnevRichterGebert1993} do not limit E1--E4, whose rank is exactly three.
Billera, Kapranov, and Sturmfels prove the sphere theorem for cellular
strings of a convex polytope ordered by a generic linear functional
(Theorem~1.2, p.~550), and their coherent-string sphere is the boundary
of the monotone path polytope (\S2, pp.~552--553)
\cite{BilleraKapranovSturmfels1994}. Their higher-dimensional projection
question is explicitly posed as the generalized Baues problem on p.~554.
None of these statements identifies the manuscript's cyclic port rule
with coherence. In particular, the local $\nu$-cube of E2 and the global
$n$-cube in Sturmfels--Ziegler's Baues projection have different roles;
identifying them would not be a type-correct comparison.
We searched the exact titles of the cited papers, ``extension space parallel
extensions punctured'', and ``oriented matroids extension space M\"obius'';
we found no source stating E3, the marked form E4, or the precise cyclic
selector and relative-cell formulas of E6 as used here. Our checked
coverage comprises the complete Sturmfels--Ziegler and cellular-strings
papers, the introduction and first section of Mn\"ev--Richter-Gebert,
and the allowable-sequence definitions of Goodman--Pollack
\cite{GoodmanPollack1980}. The general wiring statements were also
checked in the coauthor chapter, Theorems~5.1.4, 5.2.6, and 5.5.10
\cite{FelsnerGoodman2017}. This is a
report of the search, not a claim of exhaustive coverage or priority;
we supply the full deductions above.
% Owned by the polygons workstream.  Result IDs P1--P8 are stable.
\section{A.3. The image of the polygon invariant}
\label{sec:A3}

We keep the arrangement labelled.  Let $E=\{1,\ldots,n\}$, let
$\widetilde E=\{+i,-i:i\in E\}$, and assume that $\Sigma$ is a finite,
non-pencil arrangement of distinct tame pseudolines in $\mathbb {RP}^{2}$.
Multiple vertices are allowed.  The associated oriented matroid $M$ has
rank three and is simple; ``simple arrangement'' would impose the stronger
condition that $M$ is uniform.  We distinguish these two uses of
``simple.''  We choose orientations of the old pseudolines and use a
chirotope representative $\chi_M:E^3\longrightarrow\{-1,0,1\}$
calibrated with these signed half-line labels.  For example, for oriented
projective lines represented by normal vectors $a_i$, the calibration is
$\chi_M(i,j,k)=\operatorname{sgn}\det(a_i,a_j,a_k)$, and in an affine
chart the positive endpoint of the line with normal $(u_i,v_i,w_i)$ has
direction $(-v_i,u_i)$.  The geometric proofs below give the same
calibration for pseudolines.  Changing this convention requires changing
the corresponding signed labels as well; it is not an additional
ambiguity of the relative position.

An unoriented relative position of a new pseudoline $d$, distinct from
every old line, is a simple rank-three extension $N$ on $E\sqcup\{d\}$
with $N\setminus d=M$, modulo reorientation of $d$ only.  We write
$q(i,j)=\chi_N(i,j,d)$.  Thus $q$ is the chirotope of $N/d$, and
reorientation of $d$ changes $q$ to $-q$.  These correspondences use the
rank-three topological representation theorem: an oriented matroid
determines the labelled cell decomposition, and any two representations
with the same data are related by a homeomorphism.  We checked the
projective formulation in \cite[\S7.1.4, pp.~6--7]{RGZpreprint}, and the
relative uniqueness statement in \cite[Theorems~4--5, pp.~15,18--20]{BKMS2002}.
The latter fixes the old arrangement cellwise.  Consequently a
representation obtained abstractly for $N$ can be carried to the given
$\Sigma$; we do not replace the prescribed arrangement by a different
realization without this comparison.

We make explicit the projective comparison used here.  On the
spherical double covers, equal covector data identify the labelled
vertices, edges, and two-dimensional cells.  Their closures are
balls, by \cite[Theorem~2, pp.~10--11]{BKMS2002}.  Match the vertices,
extend over each pair of antipodal edges by interval homeomorphisms,
and then over each pair of antipodal two-cells by disk
homeomorphisms extending the already matched boundary circles.
On the second member of each pair use the antipodal conjugate of
the first map.  No cell is its own antipode, since its nonzero sign
vector differs from its negative.  The resulting map is an
antipodal-equivariant cellular homeomorphism of spheres, and hence
descends to a labelled homeomorphism of the projective arrangements.
This is the comparison carrying a deletion representation to the
prescribed $\Sigma$ in every use below.

\subsection{An explicit finite image criterion}

A signed boundary word is a cyclic listing $W$ of $\widetilde E$, each
label occurring once.  We may temporarily orient the cyclic listing;
reversing it is allowed when the deployed polygon is unoriented.  From
$W$ define an alternating map $q_W:E^2\longrightarrow\{-1,0,1\}$.
For $i\ne j$, put $q_W(i,j)=0$ if the two endpoint pairs
$\{+i,-i\}$ and $\{+j,-j\}$ do not alternate.  If they alternate, put
$q_W(i,j)=+1$ when $+j$ lies in the positively directed open arc from
$+i$ to $-i$, and put $q_W(i,j)=-1$ otherwise.  This is an entirely
combinatorial rule; in particular, it does not assume that $W$ is valid.

For any alternating $q:E^2\longrightarrow\{-1,0,1\}$ define
$\chi_q:(E\sqcup\{d\})^3\longrightarrow\{-1,0,1\}$ by
\begin{equation}
 \chi_q(i,j,k)=\chi_M(i,j,k),\qquad
 \chi_q(i,j,d)=q(i,j),                                      \label{eq:P-combined}
\end{equation}
and extend by alternation, with value zero on repeated entries.  The
following version of the chirotope test is explicit also in the
nonuniform case.  Its nonzero unordered triples must be the bases of a
rank-three matroid: for bases $B,C$ and $b\in B\setminus C$, some
$c\in C\setminus B$ must make $(B\setminus\{b\})\cup\{c\}$ a basis.
For every five distinct labels $\lambda,a,b,c,e$, the three signs
\begin{equation}
 \chi_q(\lambda,a,b)\chi_q(\lambda,c,e),\quad
 -\chi_q(\lambda,a,c)\chi_q(\lambda,b,e),\quad
 \chi_q(\lambda,a,e)\chi_q(\lambda,b,c)                    \label{eq:P-GP}
\end{equation}
must either include both signs or all be zero.  These are the two
conditions (CHI1)--(CHI2) in the opened source
\cite[\S7.2.3, p.~10]{RGZpreprint}.  The unsigned basis condition is
essential when zeros occur; it must not silently be omitted.

We next specify the cluster expansion used in the criterion.  Suppose
that $\chi_q$ passes this test and no row of $q$ is zero.  The rank-two
matroid of $q$ has no loops.  Its signed elements have a cyclic order of
parallel-ray classes, with opposite classes halfway round.  Choose one
of the two cyclic orientations.  In a signed ray class containing
$\varepsilon i$ and $\eta j$, choose $k$ outside the corresponding
unoriented parallel class.  Order the two labels by
\begin{equation}
 \varepsilon i \text{ precedes }\eta j
 \quad\Longleftrightarrow\quad
 -\eta\,\chi_M(i,j,k)q(i,k)=+1.                         \label{eq:P-expansion}
\end{equation}
Place the opposite class in the reversed, sign-negated order.  Denote
the resulting word by $\operatorname{Exp}_M(q)$.  The proof of P3 below
shows that the relation is a total order inside each packet and is
independent of $k$.  This gives a practical construction without real
coordinates: choose a signed representative for each class in a
half-circle of the rank-two order, sort the classes using the signs of
$q$, and sort each packet by \eqref{eq:P-expansion}.

\begin{resultenv}{P4}{Proved}{Complete image criterion}
A cyclic signed word $W$ is the deployed boundary polygon of a
non-supercritical relative position for the prescribed $\Sigma$ if and
only if the following finite tests hold, with the calibration just
specified:
\begin{enumerate}
 \item $\chi_{q_W}$ satisfies the unsigned basis axiom and
       \eqref{eq:P-GP};
 \item no row of $q_W$ is identically zero;
 \item $W$ agrees, up to a cyclic shift, with
       $\operatorname{Exp}_M(q_W)$.
\end{enumerate}
If $W$ is given without a cyclic orientation, perform these steps with
either orientation; the resulting extension is taken modulo
reorientation of $d$.  The tests include arbitrary vertex
multiplicities and any number of vertices on $d$.
\end{resultenv}

\begin{proof}
For an actual relative position, cut along a narrow M\"obius band about
$d$.  Two remaining chords cross exactly when their intersections with
$d$ were distinct.  By the Jordan separation property for a proper arc
of a disk, they cross exactly when their endpoints alternate.  In that
case their signed cyclic order gives $\chi_N(i,j,d)$, by the
rank-two chirotope convention.  Thus the pair rule gives $q_W=q$.
The combined map is the chirotope of $N$, and it passes the two axioms.
A zero row at $i$ would say that $\{i,d\}$ has rank one: $d$ would be
parallel to $i$ in the oriented matroid, and hence the same projective
pseudoline.  This is excluded.  Finally the local disk computation in
P3 is exactly the expansion of the packets, so the third condition
holds.

Conversely, the first condition produces a rank-three oriented matroid
$N$ with chirotope $\chi_{q_W}$, whose deletion of $d$ is exactly $M$.
The second condition makes $d$ a nonloop and prevents it from being
parallel to any old element.  All old elements were already distinct
and nonparallel.  Therefore the projective topological representation
of $N$ consists of the prescribed number of distinct pseudolines.
Its deletion represents $M$, so the relative uniqueness/homeomorphism
comparison stated above carries that deletion to $\Sigma$ and carries
the new pseudoline to a relative position for $\Sigma$.  P3 computes
its boundary word as $\operatorname{Exp}_M(q_W)$; the third condition
identifies that word with $W$.  This proves sufficiency and exhausts
the image.  Reversal changes the contraction orientation, and no old
label is reoriented in this comparison.
\end{proof}

There is an equivalent localization formulation that explains precisely
where Las Vergnas's theorem enters.  For independent $i,j\in E$ let
$C_{ij}$ be the signed cocircuit with entries
$C_{ij}(k)=\chi_M(i,j,k)$.  Given $q$, set
\begin{equation}
               \sigma(C_{ij})=q(i,j).                    \label{eq:P-sigma}
\end{equation}
First require that this is well defined when two ordered pairs define
the same cocircuit, and that $\sigma(-C)=-\sigma(C)$.
For every rank-one flat of $M$ contract that flat.  On the cyclic list
of $2m$ cocircuits of the resulting rank-two oriented matroid, the
values of $\sigma$ must, up to cyclic shift, be one of
\[
 0^{2m},\qquad +^{m-1}0-^{m-1}0,\qquad +^{m}-^{m}.
\]
These tests are exactly Las Vergnas's localization criterion.  We
opened its full version, including zero values, in
\cite[Definition~6.4 and Theorem~6.5, p.~25]{GP2021}; that source states
the dual version for colocalizations.  To apply it here we take $M^*$.
A subset $A\subseteq E$ has nullity two in $M^*$ precisely when
$E\setminus A$ has rank one in $M$, since
$r_{M^*}(A)=|A|-3+r_M(E\setminus A)$.  Moreover
$(M^*|A)^*=M/(E\setminus A)$, and circuits of the first matroid are
cocircuits of the second.  Thus its restrictions are exactly the
rank-two contractions above, with the same cyclic sign patterns.
This proves that the criterion applies in the direction required here,
including the nonuniform cases.  The nonzero localization determines
the extension; P2 supplies a direct uniqueness proof from the
chirotope.  The image criterion P4 adds the explicit cluster expansion
to recover the literal deployed polygon, rather than merely its
contraction.

\subsection{The proposed three local conditions are insufficient}

The phrase ``three triangles'' needs a test that can be applied before
the existence of a relative position is known.  Otherwise it already
assumes the conclusion.  We use the strongest natural local test on a
stable candidate word: an adjacent signed pair
$(\varepsilon i,\eta j)$ is a candidate empty triangle at $v_{ij}$ if,
for every $k\notin\{i,j\}$, all nonzero signs in
\begin{equation}
 \chi_M(i,j,k)q(i,j),\quad \varepsilon q(i,k),\quad
 \eta q(j,k)                                               \label{eq:P-triangle}
\end{equation}
agree.  In an actual simple arrangement these are the signs of the
vertex and the two boundary endpoints relative to the $k$th chord.
They agree exactly when that chord does not cut off the alleged
triangle.  Adjacency excludes an endpoint in its base, and testing all
other chords excludes a crossing of either side.  Hence
\eqref{eq:P-triangle} is an exact empty-triangle test for a valid word.
We apply this same finite necessary test to candidates that may fail
global compatibility.

\begin{resultenv}{P5}{Refuted}{Lead L3}
The $2n$ weight rule, the cluster-reversal rule, and the presence of at
least three empty-triangle candidates, even tested against every old
line by \eqref{eq:P-triangle}, do not characterize the image.  This
fails for a simple arrangement of six straight lines.  The complete
replacement is P4, or equivalently the localization criterion together
with the prescribed packet expansion.
\end{resultenv}

\begin{proof}
Let $a_i=(1,t_i,t_i^2)$ for $t_i=i-1$, $1\le i\le6$, and let the old
projective line $i$ have equation $a_i\cdot x=0$.  For $i<j<k$,
\[
 \det(a_i,a_j,a_k)=(t_j-t_i)(t_k-t_i)(t_k-t_j)>0.
\]
Thus all old triples have sign $+1$, all old lines are distinct, and
no three old lines are concurrent.  In particular the arrangement is
nontrivial.  Consider the cyclic word
\[
 W=(1,2,-3,-4,-6,-5,-1,-2,3,4,6,5).
\]
It is stable and exactly antipodal: the second half is the negative of
the first.  Its rank-two pair signs, in lexicographic order
$12,13,14,15,16,23,24,25,26,34,35,36,45,46,56$, are
\[
 (+,-,-,-,-,-,-,-,-,+,+,+,+,+,-).
\]
The three adjacent arcs
\[
                (2,-3),\qquad(-1,-2),\qquad(3,4)
\]
pass \eqref{eq:P-triangle} for every other old line.  For the first
arc all four sign triples are $(-,-,-)$; for the other two arcs all
four sign triples are $(+,+,+)$.  The remaining nine adjacent arcs
fail, so there are exactly three candidates.  Give each marked arc
weight one.  In their cyclic order the three complementary arcs have
weights $4,1,4$, the numbers of boundary edges they contain.  The
total weight is $1+1+1+4+1+4=12=2n$.  There are no nontrivial clusters,
so the cluster rule holds vacuously.

Nevertheless, in the combined map of \eqref{eq:P-combined}, take
$\lambda=5$, $a=1$, $b=3$, $c=6$, and $e=d$.  The three terms in
\eqref{eq:P-GP} are all $-1$.  For example they are
\[
 \chi_M(5,1,3)q(5,6)=-1,\quad
 -\chi_M(5,1,6)q(5,3)=-1,\quad
 q(5,1)\chi_M(5,3,6)=-1.
\]
This violates the chirotope axiom, so no pseudoline in the prescribed
arrangement has the candidate word.  The exact-integer program
\texttt{polygons\_l3\_search.py} checks all determinants, all twelve
triangle tests, the stated violation, and a failure of the full signed
basis-exchange axiom.  Its complete output is
\texttt{audit/polygons\_l3\_search\_output.json}.  A second direct
verifier, \texttt{polygons\_l3\_verify.py}, repeats the certificate
without searching.  Thus the counterexample concerns the prescribed
$M$, not a polygon borrowed from a different arrangement.
\end{proof}

The search also tested every stable signed word for the four- and
five-element moment-curve arrangements (48 and 384 oriented words),
and all 3840 such words for the six-line test arrangement of A.1;
none of those searches produced this obstruction.  The six-element
moment-curve example was found at candidate 49.  These finite negative
searches do not assert minimality among all arrangements.

\subsection{Recovering the weighted graph}

\begin{resultenv}{P6}{Proved}{Reconstruction of vertices, arrows, rotations, and weights}
For a fixed labelled nontrivial arrangement $\Sigma$, its full image of
deployed polygons determines all non-supercritical relative positions,
all ordinary specialization arrows, their directions and weights, and
their cyclic incidences at each position.  Adjoining the known $n$
supercritical positions recovers the corridor arrows and their
weights.  Together with the prescribed transition transport used in
constructing the surface, these data determine its cellularly embedded
weighted graph, up to label-preserving isomorphism of cellular maps.
The reconstructed graph alone does not assert that every face has
order four; that is the separate local-link assertion of A.1.
\end{resultenv}

\begin{proof}
P2 below reconstructs the extension at every polygon.  Its signed
cocircuit localization is \eqref{eq:P-sigma}.  A projective old vertex
$v$ is a pair $\{C_v,-C_v\}$ of cocircuits; write
$F_v=C_v^0\subset E$ for the set of old lines through it, and
$m_v=|F_v|$.  An ordinary immediate specialization is obtained by
replacing the two values at $\pm C_v$ by zero and leaving every other
localization value unchanged, provided the resulting localization
passes the extension criterion and is not supercritical.  This rule
is finite and orientation-independent: simultaneous negation of all
values just reorients $d$.

We verify the geometric equivalence of the rule.  An elementary sweep
to $v$ moves no other old vertex across $d$, so it has exactly this
effect on the localization.  Conversely, if the rule yields an
extension $N'$, represent it with its new line through $v$.  Choose a
disk about $v$ containing no other old vertex.  Resolving the new line
to either side inside that disk gives the two localizations obtained
by replacing the zero at $C_v$ by the two possible signs; all other
coordinates are unchanged.  Exactly one has the given nonzero sign
of the original extension $N$.  Uniqueness of the extension therefore
identifies that local resolution with $N$.  This comparison is a
homeomorphism preserving the old arrangement cellwise, so the local
resolution proves that the required empty fan exists for the original
relative position.  In a cut-open disk it is precisely the
specialization triangle, subdivided by the $m_v$ concurrent old
lines.  There are $m_v-1$ boundary edges in its base, which proves
that the arrow weight is $m_v-1$.

At $N'$, the class $F_v$ is visible directly from the endpoint
nonalternation relation.  Its two boundary packets give the two
generization intervals.  Reversing one packet performs the
corresponding resolution; the other packet stays fixed.  At $N$, the
same marked interval is the specialization base.  Thus the signed
boundary word specifies the interval occupied by each incident
ordinary arrow, not only its two end vertices.  Their circular order
is the required local rotation, and the lengths of the complementary
intervals give all corner weights.  The intervals partition the
$2n$ boundary edges between transition bases and corner intervals,
which also proves the $2n$ weight identity.

For a supercritical position $i$, use the parallel extension
$N_i$ with $d$ parallel to $i$.  Its localization is
$\sigma_i(C)=C(i)$, and its contraction entries are
$q_i(j,k)=\chi_M(j,k,i)$.  Let $N$ be substable, with its unique old
vertex $v$ on line $i$.  There is a corridor arrow $N\to N_i$ exactly
when, after choosing the sign of the new element,
\[
          \sigma_i(C)\ne0 \quad\Longrightarrow\quad
          \sigma_N(C)=\sigma_i(C) .                     \label{eq:P-corridor}
\]
Necessity is immediate by collapsing the corridor.  For sufficiency,
contract $i$.  The cyclic rank-two localization criterion shows that,
with the zero at $v$ prescribed, there are just the two choices
corresponding to the two local sides of $i$ at $v$.  A narrow strip
about $i$, with the new line crossing $i$ at $v$ and following either
side elsewhere, realizes these choices.  Its localization away from
$i$ is $C(i)$; on $i$ it is the prescribed rank-two pattern.
Therefore \eqref{eq:P-corridor} and uniqueness identify $N$ with one
of these narrow-strip positions.  This is the required corridor.
The $n-m_v$ lines not through $v$ cross the corridor as disjoint
transversal chords, dividing its side into $n-m_v+1$ edges.  Hence
the ordinary corridor weight is $n-m_v+1$; its modified weight at
the supercritical vertex is $m_v-1$, and their sum is $n$.
The order of the corridor flags at $i$ is the order of the old
vertices on the twofold cover of $i$, also determined by $M/i$.

Finally, a graph with cyclic orders at its vertices and specified
identifications of the two local orientation choices across its edges
is a signed rotation system.  Thicken each vertex to a disk and each
edge to a band with that identification.  Its boundary cycles are
determined by following the finite permutations of flags; attaching
one disk to each such cycle gives the unique cellular map with these
data.  Transition transport is part of the construction being
reconstructed, and its flag identifications are obtained from the
labelled local sweep or the corridor.  This establishes the asserted
reconstruction without deducing quadrangularity or embeddability of
the original large disks from the mere existence of the signed
rotation system.
\end{proof}

\section{A.4. Proofs of the disk criterion and the two polygon assertions}
\label{sec:A4}

\subsection{The precise disk criterion}

\begin{resultenv}{P1}{Proved}{Sufficiency in R1}
Let $\Delta$ be a closed disk, with finitely many labelled tame proper
arcs $\delta_i$ whose $2n$ endpoints on $\partial\Delta$ are distinct.
Suppose each pair meets at most once, at an interior transverse
crossing, and disjointness together with equality is an equivalence
relation on the arcs.  If there are at least two equivalence classes,
this system is the cut-open system of a projective pseudoline
arrangement and an added pseudoline.  If there is just one class,
the same assertion holds exactly when its endpoints form two
consecutive packets, each containing one endpoint of each chord.
The reconstructed old arrangement is nontrivial exactly when its
pairwise intersection points are not all the same point.
\end{resultenv}

\begin{proof}
We first prove that the packet condition is automatic with at least
two classes.  Fix a class $C$ and an arc $\beta$ in another class.
Every arc in $C$ crosses $\beta$ once.  The pairwise disjoint arcs of
$C$ therefore separate the two endpoints of $\beta$ and occur in a
linear chain along it.  To see explicitly that there is no branch,
cut $\Delta$ along $\beta$.  In either resulting half-disk the pieces
of arcs in $C$ are pairwise disjoint arcs joining distinct points of
the cut to the opposite boundary; noncrossing forces the same linear
order, reversed when the boundary is traversed in the opposite
direction.  Thus the endpoints of $C$ occur in two packets with
opposite orders.

No endpoint of an arc outside $C$ can lie strictly between two
consecutive endpoints in either packet.  Indeed those two arcs of
$C$ bound a middle strip.  An arc starting in that strip and crossing
each of its two boundary arcs exactly once cannot end in a component
outside both: after the first crossing it is outside one boundary,
and to reach the other boundary it would have to recross the first.
This contradicts the one-crossing hypothesis.  Since every outside
arc crosses both boundaries, neither endpoint can start in the
middle strip.  It follows that the two packets are consecutive among
all chord endpoints, not merely among endpoints belonging to $C$.

Collapse each packet to a point of the boundary.  The quotient of the
disk is again a disk.  For completeness, one may describe the quotient
map in polar coordinates: choose a nondecreasing degree-one boundary
map $\phi$ that is constant exactly on the packet intervals, and use
$t\mapsto (1-r)t+r\phi(t)$ on the circle of radius $r<1$.  These maps
are strictly increasing for $r<1$, extend continuously to the boundary
collapse, and induce a continuous bijection from the quotient onto a
disk.  Compactness gives a homeomorphism.  Finite piecewise linear
choices give the corresponding tame local model.

Choose a representative chord from each of the $r$ classes.  They
cross pairwise.  Their $2r$ boundary points have order
$C_1,\ldots,C_r,C_1,\ldots,C_r$: after cutting at the two endpoints
of $C_1$, every other chord has one endpoint in either half, and two
such chords cross exactly when their orders in the two halves agree.
This proves the displayed order.  Consequently the involution pairing
the two points of each class is an order-preserving, fixed-point-free
boundary involution.  Identify this boundary by that involution.
After a boundary homeomorphism it is the usual antipodal
identification, so the quotient is $\mathbb {RP}^2$.

The image of the boundary is a new pseudoline $d$.  Each old chord
closes to a simple curve, with its two endpoints identified and no
other self-identification.  It crosses $d$ once, hence is an essential
one-sided simple curve and is a pseudoline.  Chords from different
classes retain their unique interior transverse intersection.  Chords
in the same class meet only at the boundary vertex created for that
class.  Their opposite packet orders give, in the local disk after
gluing, the cyclic order of rays of a pencil: each pair of old curves
crosses there, and each crosses $d$.  Thus all required pairwise
intersection and crossing conditions hold.  Expanding the packet
vertices inside a narrow band about $d$ recovers the original disk
system, up to a labelled homeomorphism.

For one class the same construction works when two packets exist.
Conversely, if every old chord is disjoint, every pair of old
pseudolines must have met on $d$.  Since each old pseudoline meets
$d$ only once, all those intersections are a single vertex, and a
small neighborhood of that vertex gives exactly two packets.  The
condition is necessary.  Finally the quotient construction lists all
old intersection points, interior crossings and class vertices;
their being a single common point is precisely the pencil case.
\end{proof}

The packet condition occurs explicitly in the transcription of
154/035, after the disjointness-equivalence clause.  The transcription
marks some bracketed wording about antipodality and uniqueness as
doubtful or deleted.  P1 does not depend on either reading of that
bracket: it proves the packet assertion where it is automatic, and
states its exact necessity in the remaining case.  If R1 is read as
asserting sufficiency of the two abbreviated conditions for all finite
systems without this qualification, it is false: three disjoint small
chords cutting off three separate boundary caps have word
$(1,-1,2,-2,3,-3)$, and admit no two packets.  The finite endpoint
certificate is in \texttt{polygons\_certificates.py}.  This degenerate
example does not arise from the nontrivial rank-three arrangements in
the question.  For $n=0$ the empty disk system may be closed with an
arbitrary added line; for $n=1,2$ every old arrangement is a pencil.
Those cases are outside the nontrivial rank-three hypothesis.  No
assertion about infinite or nontame systems is being used.

\subsection{Injectivity and cluster reversal}

\begin{resultenv}{P2}{Proved}{R2: the literal polygon determines the relative position}
For the fixed labelled nontrivial arrangement $\Sigma$, the map from
non-supercritical relative positions to literal deployed signed
polygons is injective.  The polygon determines the rank-two
contraction by the signed pair-alternation rule, including all its
parallel classes.  The only ambiguity is reorientation of the added
element, corresponding to reversal of the chosen cyclic orientation.
\end{resultenv}

\begin{proof}
The Jordan separation argument in P4 recovers every value
$q(i,j)=\chi_N(i,j,d)$ from the word.  In particular, no guess at a
cluster decomposition is needed: $i$ and $j$ are in the same cluster
if and only if their endpoint pairs do not alternate.  For an actual
word this is exactly equality of their intersections with $d$, hence
is an equivalence relation.  We know $\chi_N$ on every old triple
from $M$, and on every triple containing $d$ from $q$.  Alternation
and zeros on repetitions determine all remaining ordered triples.
Thus \eqref{eq:P-combined} determines the entire extension.  The
relative topological representation comparison identifies its
relative position in the fixed arrangement.  Reversing $q$ changes
exactly the chirotope entries containing $d$, which is precisely
reorientation of $d$.

Equivalently, $q$ gives the cocircuit localization by
\eqref{eq:P-sigma}, and Las Vergnas's uniqueness of a single-element
extension with a prescribed localization gives the same conclusion.
The direct argument above identifies the exact data involved and
does not need to invoke the stronger existence theorem.  In
particular it avoids confusing the contraction with the deployed
polygon: the latter has the additional within-packet orders dealt
with in P3.
\end{proof}

\begin{resultenv}{P3}{Proved}{R3: the exact antipodal permutation}
If the vertices of $\Sigma$ on $d$ have line clusters of sizes
$k_1,\ldots,k_\nu$, the deployed word consists of paired consecutive
packets in opposite order.  Let $H$ be the half-turn on the $2n$
positions and let $R$ reverse the order separately within every
packet, fixing singleton packets.  The label-negation involution is
$H\circ R=R\circ H$.  It equals $H$ exactly when $d$ is stable.
The within-packet order is \eqref{eq:P-expansion}.
\end{resultenv}

\begin{proof}
At an old vertex $v$ on $d$, choose a small disk containing no other
old vertex.  The $k$ old line germs and the germ of $d$ form a pencil
of transverse proper arcs.  A homeomorphism of this disk carries
the finite star to straight radial segments, preserving its cyclic
order.  Take $d$ locally to be $y=0$ and write the other lines as
$x=b_i y$, with distinct $b_i$.  The two nearby boundary arcs meet
them at $x=\epsilon b_i$ and $x=-\epsilon b_i$.  Their orders are
reversed.  This proof is topological: the straightening is only of
one local star, and is available whether the whole arrangement is
stretchable or not.

After the packets are collapsed, the endpoints are the rank-two
cyclic order on $d$.  Corresponding packets are opposite and have
the same cardinality.  Hence the full word has the form
\[
       B_1\cdots B_r\,(-B_1)^{\rm rev}\cdots(-B_r)^{\rm rev}.
\]
Its first and second halves each have $n$ entries.  The formula for
label negation follows entry by entry.  If a packet has at least two
entries, its reversal is nonidentity, so label negation is not the
half-turn.  If every packet is a singleton, the reversals are all
identities.  Thus the equivalence with stability is exact.

We also verify the signed expansion formula, needed for the
algorithm.  It involves only two lines $i,j$ of one cluster, $d$,
and an old line $k$ outside that cluster.  This pencil-with-transversal
configuration is straightenable by choosing an affine chart with
$d$ at infinity.  Write its normal rows as $a_i=(w_i,z_i)$,
$a_j=(w_j,z_j)$, $a_k=(w_k,z_k)$ with $w_\ell\in\mathbb R^2$.
After reorienting $i,j$ by $\varepsilon,\eta$, their vectors
$\varepsilon w_i$ and $\eta w_j$ lie on the same ray. Write
$w_j=\lambda w_i$, where $\lambda\ne0$ and
$\operatorname{sgn}\lambda=\varepsilon\eta$. Elementary expansion gives
\[
 \det(a_i,a_j,a_k)=-(z_j-\lambda z_i)\det(w_i,w_k).
\]
On a large circle of radius $R$, the signed endpoints are near
$J(\varepsilon w_i)/\|w_i\|=J(\eta w_j)/\|w_j\|$, where $J$
rotates by $\pi/2$. Their angular offsets are ordered by
$\varepsilon z_i/\|w_i\|$ and $\eta z_j/\|w_j\|$.
Consequently the sign of the second offset minus the first is the sign of
$-\eta\chi_M(i,j,k)q(i,k)$, as in
\eqref{eq:P-expansion}.  The order is geometric and independent of
the choice of transversal $k$, so the formula is independent of $k$
and is transitive for the whole packet.  Only the local pencil and
transversal were straightened, which proves the formula for arbitrary
pseudoline arrangements.

A smallest nontrivial explicit failure of the half-turn claim is
given by old equations $x=0$, $x=z$, $y=0$ and $d:z=0$.  The old
normal determinant is $+1$, so the old arrangement is nontrivial.
The word is $(2,1,3,-1,-2,-3)$; its half-turn takes $2$ to $-1$.
The exact computation in \texttt{polygons\_certificates.py} verifies
the arrangement, the word, and the full chirotope axiom.  Therefore
the literal antipodality assertion on 154/018 is false for
nonstable positions; the correction R3 is true.
\end{proof}

\subsection{A local fact used in the surface analysis}

\begin{resultenv}{P7}{Proved}{Commuting resolutions at distinct vertices}
Let $d$ be a non-supercritical relative position through distinct old
vertices $v_1,\ldots,v_\nu$.  Resolving any chosen subset to either
local side can be performed independently.  The resulting
generizations are indexed by $\{0,+,-\}^{\nu}$, where zero means
unresolved, and the $2^\nu$ fully resolved positions are distinct.
In particular two distinct chosen vertices give four distinct
positions in each ordinary commuting square.
\end{resultenv}

\begin{proof}
Distinct vertices on $d$ cannot share an old line: that old line
would meet $d$ twice unless it were $d$, contrary to the
non-supercritical assumption.  Choose mutually disjoint small disks
about the $v_i$, meeting no other old vertices.  In each disk replace
the arc of $d$ by either of the two nearby transverse arcs, or leave
it unchanged.  Outside these disks leave $d$ fixed.  The resulting
curve stays simple and meets each old line once, because an old line
is affected in at most one of the disks and the local replacement
retains its unique crossing.  Operations in disjoint disks commute.
The localization changes at exactly the corresponding pair of old
cocircuits, from zero to the chosen sign.  Different index vectors
therefore give different localizations and different relative
positions, even after forgetting the orientation of $d$: an old
vertex not on $d$ fixes the relative global sign, or, equivalently,
one can use an unchanged nonzero entry of the contraction between two
different clusters.  Such an entry exists because $M$ has rank three
and $d$ is non-supercritical.  This proves the asserted cube and the
distinctness of its square vertices.
\end{proof}

P7 proves the ordinary commuting part of a corner argument, but does
not itself identify every corner of the disk-glued space with such a
square.  In particular two specialization bases can share an old
line, so their future apices cannot both lie on a distinct added
pseudoline; that case must be treated by the corridor/supercritical
construction.  We retain this distinction to avoid using a local
cube as an unproved global gluing theorem.

\begin{resultenv}{P8}{Proved}{The exposed-fan or corridor alternative}
Suppose that a non-supercritical new pseudoline $d$ passes through
exactly one old vertex $v$.  On the spherical double cover, let $\alpha$ be either
open semicircle of $d$ between the two lifts of $v$.  Choose an old
line $i$ through $v$ adjacent to $d$ on one side, and let $L$ be the
closed lune bounded by $\overline\alpha$ and the corresponding
semicircle of $i$.  If $L$ contains an old vertex in its interior,
then an empty specialization fan has its base in $\alpha$.  If it
contains no old interior vertex, $d$ is in the corridor position next
to $i$ on that side.  In particular, a half-circle between its
generization ports contains an
ordinary specialization or is an exceptional corridor half-circle.
\end{resultenv}

\begin{proof}
Here ``adjacent'' refers to the cyclic order of the old line germs
and the germ of $d$ at $v$.  An old line through $v$, other than $i$,
cannot enter the interior of $L$: it has no further intersection
with either boundary line, and its germ is outside this sector.
Every old line not through $v$ meets each wall once and cuts $L$ in
one proper arc joining its $d$ wall to its $i$ wall.  Indeed each
wall semicircle contains exactly one of the two antipodal
intersections with that old line, and these are the only boundary
intersections.  We call these arcs the wires and direct them from
the $d$ wall to the $i$ wall.

We justify that the resulting finite directed graph of old vertices
has no directed cycle.  Let $N$ be the oriented matroid of the whole
arrangement including $d$, use $g=d$ as the distinguished element
at infinity, and use $f=i$ as objective.  Thus the affine space in
the terminology of \cite[pp.~30--31]{SZ1993} is the hemisphere
$S_d^+$ containing $L$.  We choose the sign of $i$ so that the
program orientation points toward the $i$ wall in this lune.  The
definition of its graph $G_i$ assigns a direction to an edge along
a wire by the sign of $i$ at the wire's point at infinity on $d$.
That sign is constant on the chosen open $d$ wall, and is nonzero:
a traversing wire does not pass through $v=d\cap i$.  Hence every
wire edge is directed, and its direction is precisely the one just
specified.  The interior graph in $L$ is a subgraph of $G_i$.

The program satisfies the hypotheses of the cited definition.  In
particular $d$ is not a loop and $i$ is not a coloop.  For the latter,
another old line through $v$, an old line not through $v$, and $d$
form a rank-three basis after deleting $i$; such a line not through
$v$ exists because the old arrangement is not a pencil.  The proof
of \cite[Corollary~4.5, p.~40]{SZ1993} shows that no rank-three
oriented matroid program is non-Euclidean.  By its Definition~3.8(i)
and Proposition~3.9 (p.~33), such a program has no directed cycle
(a path may use undirected edges, but a directed cycle must contain
at least one directed edge).  In our subgraph every edge is directed,
so it is an ordinary finite directed acyclic graph.

If there is an interior vertex, choose a source $s$ of this graph.
For each old line through $s$, its incoming arc from the $d$ wall
has no old vertex before $s$; otherwise the last such vertex and
the intervening edge would contradict the source property.  These
incoming arcs have distinct endpoints on $\alpha$ and are mutually
disjoint except at $s$.  They form a star that partitions $L$ into
regions.  Exactly one region has the $i$ wall on its boundary.  Each
outgoing branch at $s$ must lie in that region: if it initially lay
in any other region, then to reach its endpoint on the $i$ wall it
would have to cross an incoming arc again or meet the $d$ wall
again.  Either event contradicts the unique-intersection property
of the pseudolines.  Thus the incoming branches are consecutive in
the cyclic order at $s$.  Their two extreme branches and the
interval of $\alpha$ between their endpoints bound a fan disk $F$.

No old line not incident to $s$ enters the interior of $F$.  Such a
line cannot cross either extreme incoming arc, since an intersection
there would be a vertex preceding $s$.  If it entered through the
base on $d$, it would have to leave through one of those extreme
arcs or through $d$ a second time; it cannot leave through $s$
because it is not incident to $s$.  These alternatives are all
impossible.  The intermediate incoming arcs only subdivide $F$
into the usual empty triangles.  We can therefore move the base
arc of $d$ through this empty fan to $s$, without changing any
other old intersection.  This gives the required ordinary
specialization and places its full base in $\alpha$.

If there is no old interior vertex, all wires in the open lune are
pairwise disjoint.  Move the $d$ wall to a curve just inside the
$i$ wall, keeping its endpoints at the two lifts of $v$.  This move
can be performed without passing an old vertex: the disjoint
wire arcs cut the lune into finitely many disk strips, with fans
allowed at their endpoints on $i$, and within each such strip we
slide the boundary interval across a rectangle.  The slides agree
on their common wire sides.  Equivalently they give an isotopy of
the new arc through transverse arcs, preserving one intersection
with each wire.  It extends over the complementary cells to a
homeomorphism preserving every labelled old line.  The final arc
is a narrow push-off of $i$ crossing it at $v$, on the specified
local side.  This is exactly the corridor position.  Conversely,
such a corridor lune contains no old interior vertex.  The two
alternatives prove the statement and justify the exceptional case
described on 154/093 independently of its doubtful marginal text.
\end{proof}
\section{Lead L4: the two kinds of arc data}

\begin{resultenv}{T4}{Refuted}{The 1986 tethers are not the proper chords of R1}
The literal identification in L4 is false. The additional arc datum
described for an interior lieu on cote 156-4/032 has an interior end
and an end on the boundary. A chord of R1 has both ends on the boundary.
These are different types of embedded arc, even in an ordinary disk.
The exact elementary relationship is the following: a proper chord
together with a marked interior point splits into two boundary tethers
with a common interior end and disjoint interiors; conversely, two
such tethers with distinct boundary ends glue uniquely to a proper
chord. For families, the pairwise single-crossing and parallel-class
conditions of R1 must still be imposed separately.
\end{resultenv}
\begin{proof}
We first state exactly what reading this uses. On the original scan
156-4/032, in the paragraph dated 12 June, the additional datum is an
equivalence class of arcs joining an interior lieu to $\partial D$.
We inspected the whole scan and an enlarged crop. Michel Hua's project
was consulted as a secondary reading witness; its corresponding
paragraph leaves the additional datum partly unread, so we do not
take the identification of that datum from the project. On 154/035
the proper-chord requirement is explicit:
$\partial d=d\cap\partial\Delta$. Our assertion compares these typed
arc data. It makes no claim that the fragmentary 1986 discussion is
a complete independent axiom system for pseudo-disks.

Here is a complete counterexample to identifying the two kinds of
arc data, rather than a purported model of all unspecified 1986 axioms.
Let $D=[-1,1]^2$, let the interior lieu be the singleton
$A=\{(0,0)\}$, and take
\[
 \gamma:[0,1]\longrightarrow D,\qquad\gamma(t)=(t,0).
\]
The set $A$ is compact and contractible, and $\gamma$ is an embedded
arc from $A$ to $\partial D$, meeting the boundary only at its final
endpoint. It therefore supplies the boundary-joining arc datum
under discussion. But
\[
 \partial\gamma=\{(0,0),(1,0)\},\qquad
 \gamma\cap\partial D=\{(1,0)\},
\]
so it is not a proper chord. By contrast,
$c(t)=(2t-1,0)$ is a proper chord of the same disk and has two
boundary endpoints. The number of arc endpoints in the ambient
boundary is invariant under every homeomorphism of pairs
$(D,\partial D)\to(D',\partial D')$. Thus no such homeomorphism
identifies these arc data. The example has a finite polygonal
description; \texttt{l4\_endpoint\_certificate.py} checks its endpoint
incidences exactly, without floating-point decisions.

For the positive statement, a homeomorphism $[0,1]\to c$ and a marked
point corresponding to $t_0\in(0,1)$ split $c$ into the two intervals
$[0,t_0]$ and $[t_0,1]$. Each has one interior and one boundary end,
and their intersection is precisely the marked point. Conversely the
union of two embedded intervals with exactly that common endpoint,
otherwise disjoint, is homeomorphic to $[0,1]$. If their other ends
are distinct points of $\partial D$ and their remaining points lie
in the interior, this union is a proper chord. These operations
are inverse on the stated data, up to homeomorphism preserving the
marked point and the unordered pair of chosen tethers; ordering the
pair instead corresponds to orienting the chord. Without the marked point,
pairing, and second tether there is no inverse operation: for example
the single $\gamma$ above can be completed to a chord ending at
$(-1,0)$ or at $(0,1)$, with different boundary endpoint data.
Nothing in this operation enforces a restriction on intersections
of different completed chords. R1's additional conditions are
therefore exactly extra hypotheses, not consequences of the tether
description.

The sentence about two interior lieux on 156-4/045 has a separate
reading problem: the visible affirmative clause about a subdivision
conflicts with the following conclusion, which suggests a missing
negation. We retain both readings. Under the affirmative reading
the two lieux have a common subdivision; under the negative reading
they illustrate an additional compatibility obstruction. Neither
reading changes the endpoint type on /032, the counterexample, or
the cutting-and-gluing proof. Hence no mathematical conclusion here
depends on resolving that doubtful negation. The suggested historical
influence between the two folders is not inferred from the repeated
word for an arc.
\end{proof}
\section{Checked scope and statements not used}

The numbered results constitute the checked part of this paper. The
accompanying results register and check-event table distinguish a
targeted proof check, an independent re-proof, an exact finite
computation, and an earlier result rechecked when used. No numerical
experiment is substituted for the general local-link proof.

We do not assert a complete formalization in a proof assistant. We do
not infer a mathematical identification from a shared French term in
two manuscript folders. We do not assert a priority claim for the
deductions here. The title and source searches used to look for prior
versions are recorded in the report and source notes.

There are no additional unchecked numbered mathematical assertions.
The following bibliographical or interpretive points remain separate
from the proved statements. The full 1999 book chapters and the 1978
Las Vergnas original were not accessible in this run; their entries
are marked not opened, and the existence criterion is derived from
the opened full formulation of Galashin and Postnikov. The doubtful
negation on cote 156-4/045 is retained in both readings in T4, with
the same endpoint-type conclusion under each. A possible historical
influence between the folders is not decided. Neither issue is used
as a premise in a proof.

The finite scope excludes infinite and non-locally-finite families,
and the rank-three scope excludes empty arrangements and pencils.
P1 treats the degenerate disk question separately. The questions on
tame spaces in cotes 120--121, 153, and 155 belong to the other
research request and enter no assertion here.

\section{Reproduction and verification record}
The archive contains all mathematical source programs, accepted
localization lists, exact integer inputs, complete boundary words,
edge endpoint permutations, local-link checks, orientation certificates,
and source-reading notes. The main program sequence is
\texttt{counts\_extensions.py},
\texttt{counts\_independent\_check.py},
\texttt{counts\_gluing.py}, and
\texttt{extension\_compare\_cells.py}. Separate scripts certify the
fixed-straight count, the polygon counterexamples, and the tether
endpoint obstruction. The report supplies their scopes, runtimes,
source locators, and the PDF's typesetting audit. The mathematical
proofs are entirely in this paper; the scripts make the finite
certificates independently reproducible.

\nocite{*}
\bibliographystyle{apalike}
\bibliography{RESEARCH_SA11_references}
\end{document}
